\documentclass[11pt,a4paper]{article}

\usepackage[utf8]{inputenc}
\usepackage[T1]{fontenc}
\usepackage[spanish,es-nodecimaldot]{babel}
\usepackage{lmodern}
\usepackage[a4paper,margin=2.2cm]{geometry}
\usepackage{xcolor}
\usepackage{hyperref}
\usepackage{enumitem}
\usepackage{booktabs}
\usepackage{longtable}
\usepackage{tabularx}
\usepackage{array}
\usepackage{amsmath}
\usepackage{amssymb}
\usepackage{listings}
\usepackage[most]{tcolorbox}
\usepackage{fancyhdr}
\usepackage{microtype}
\usepackage[strings]{underscore}
\usepackage{textcomp}

% -----------------------------------------------------------------------------
% PARÁMETROS EDITABLES
% Cambia solamente estos valores para personalizar la guía.
% -----------------------------------------------------------------------------
\newcommand{\NombreEstudiante}{Nombre: \rule{6cm}{0.4pt}}
\newcommand{\FechaPrueba}{Por confirmar}
\newcommand{\MinutosSesion}{40}
\newcommand{\MetaSemanal}{5}

% true  = guía completa con pauta y solucionario.
% false = versión de práctica sin respuestas (ideal para autoevaluarse).
\newif\ifmostrarSoluciones
\mostrarSolucionestrue

\definecolor{azulUCM}{HTML}{17365D}
\definecolor{azulClaro}{HTML}{EAF2F8}
\definecolor{verdeClaro}{HTML}{E9F7EF}
\definecolor{amarilloClaro}{HTML}{FFF8DC}
\definecolor{rojoClaro}{HTML}{FDEDEC}
\definecolor{grisCodigo}{HTML}{F4F6F7}
\definecolor{sqlKeyword}{HTML}{005CC5}
\definecolor{sqlString}{HTML}{A31515}
\definecolor{sqlComment}{HTML}{22863A}

\hypersetup{
  colorlinks=true,
  linkcolor=azulUCM,
  urlcolor=blue,
  pdftitle={Guía de lectura activa -- Unidad 1 de Bases de Datos},
  pdfauthor={Material de estudio basado en INF-324 y documentación oficial de Oracle}
}

\pagestyle{fancy}
\fancyhf{}
\lhead{INF-324 -- Unidad 1}
\rhead{Oracle SQL}
\cfoot{\thepage}
\setlength{\headheight}{14pt}

\lstdefinelanguage{OracleSQL}{
  morekeywords={SELECT,FROM,WHERE,DISTINCT,AS,AND,OR,NOT,BETWEEN,IN,LIKE,IS,NULL,
    ORDER,BY,ASC,DESC,GROUP,HAVING,JOIN,INNER,LEFT,RIGHT,FULL,OUTER,CROSS,ON,USING,
    CASE,WHEN,THEN,ELSE,END,COUNT,SUM,AVG,MIN,MAX,ROUND,TRUNC,MOD,UPPER,LOWER,
    INITCAP,CONCAT,SUBSTR,LENGTH,INSTR,TO_CHAR,TO_NUMBER,TO_DATE,NVL,NVL2,
    COALESCE,NULLIF,DECODE,SYSDATE,MONTHS_BETWEEN,ADD_MONTHS,NEXT_DAY,LAST_DAY},
  sensitive=false,
  morecomment=[l]{--},
  morestring=[b]',
}

\lstset{
  language=OracleSQL,
  backgroundcolor=\color{grisCodigo},
  basicstyle=\ttfamily\small,
  keywordstyle=\color{sqlKeyword}\bfseries,
  stringstyle=\color{sqlString},
  commentstyle=\color{sqlComment}\itshape,
  numbers=left,
  numberstyle=\tiny\color{gray},
  stepnumber=1,
  numbersep=8pt,
  frame=single,
  rulecolor=\color{gray!35},
  breaklines=true,
  showstringspaces=false,
  upquote=true,
  tabsize=2,
  columns=fullflexible
}

\newtcolorbox{objetivo}{
  colback=azulClaro,
  colframe=azulUCM,
  title=Objetivo de lectura,
  fonttitle=\bfseries
}
\newtcolorbox{pausa}{
  colback=amarilloClaro,
  colframe=orange!70!black,
  title=Pausa activa,
  fonttitle=\bfseries
}
\newtcolorbox{clave}{
  colback=verdeClaro,
  colframe=green!45!black,
  title=Idea clave,
  fonttitle=\bfseries
}
\newtcolorbox{errorbox}{
  colback=rojoClaro,
  colframe=red!60!black,
  title=Error frecuente,
  fonttitle=\bfseries
}
\newtcolorbox{prueba}{
  colback=violet!5,
  colframe=violet!70!black,
  title=Conexión con la primera prueba,
  fonttitle=\bfseries
}
\newtcolorbox{rutina}{
  colback=blue!3,
  colframe=azulUCM,
  title=Rutina diaria de estudio,
  fonttitle=\bfseries
}
\newtcolorbox{plantilla}[1]{
  enhanced,
  breakable,
  colback=gray!3,
  colframe=azulUCM!80,
  title={Plantilla textual: #1},
  fonttitle=\bfseries
}

\newcommand{\campo}[1]{\texttt{#1}}
\newcommand{\tabla}[1]{\texttt{#1}}
\newcommand{\nivel}[1]{\textbf{#1}}
\newcommand{\casilla}{\(\square\)}

\title{\textbf{Guía de lectura activa de la Unidad 1}\\[0.3em]
\large Consultas Oracle SQL sobre el esquema OEHR\\[0.5em]
\normalsize Preparación progresiva para la primera prueba de INF-324}
\author{Guía construida desde las clases U1 L1--L6, la pauta de prueba y Oracle Database 21c}
\date{Septiembre de 2026}

\begin{document}

\maketitle
\thispagestyle{empty}

\begin{center}
\NombreEstudiante \hfill \textbf{Fecha de prueba:} \FechaPrueba
\end{center}

\begin{objetivo}
Al terminar esta guía deberías poder traducir un requerimiento escrito en español a una consulta
Oracle SQL correcta, justificar cada cláusula y detectar errores frecuentes antes de comparar tu
respuesta con la pauta. El foco es comprender, no memorizar consultas completas.
\end{objetivo}

\vfill

\begin{center}
\fbox{\parbox{0.9\textwidth}{
\textbf{Cómo usar esta guía.} Lee una sección corta, tapa el ejemplo, intenta reconstruirlo y recién
después compáralo. Al llegar a las preguntas, escribe tus respuestas antes de mirar el solucionario.
Cada módulo contiene \textbf{3 preguntas muy fáciles, 2 intermedias y 1 difícil}. Si fallas una
pregunta, vuelve al párrafo que explica esa decisión, no solamente a la solución.
}}
\end{center}

\begin{prueba}
La evaluación histórica usada como referencia pondera un \textbf{30\%}, dispone de
\textbf{90 minutos}, contiene \textbf{6 requerimientos} y usa una pauta total de
\textbf{35 puntos}: 5 por formalidad y 5 por cada requerimiento. Estos datos sirven para practicar;
confirma con el docente la fecha, duración y ponderación de la evaluación vigente.
\end{prueba}

\vfill

\begin{flushleft}
\textbf{Base principal del contenido local:}
\begin{itemize}
  \item scripts \campo{leccion-01.sql} a \campo{leccion-06.sql};
  \item presentaciones \campo{leccion-01.ppt} a \campo{leccion-06.ppt};
  \item evaluación \campo{es1-consultas-sql.docx};
  \item pauta revisada \campo{pauta-es1-consultas-sql.sql}.
\end{itemize}
\end{flushleft}

\newpage
\tableofcontents
\newpage

\section{Mapa de la unidad y estrategia de resolución}

La Unidad 1 avanza desde una consulta de una sola tabla hasta consultas que combinan tablas,
agrupan filas y emplean subconsultas. La evaluación de referencia contiene seis requerimientos y la
pauta SQL revisada propone una solución para cada uno. Por eso esta guía distingue entre:

\begin{itemize}
  \item \textbf{contenido confirmado por la pauta}: concatenación, filtros con patrones, funciones
  de grupo, \campo{GROUP BY}, \campo{HAVING}, uniones, auto-unión y subconsultas;
  \item \textbf{herramientas necesarias enseñadas en clases}: consulta básica, operadores,
  funciones de texto, número y fecha, conversiones, nulos y ordenamiento;
  \item \textbf{ampliaciones presentes en las diapositivas}: \campo{CASE}, \campo{DECODE},
  \campo{COALESCE}, distintos tipos de \campo{JOIN} y funciones adicionales de fecha.
\end{itemize}

\subsection{Ciclo semanal para leer y practicar todos los días}

Repite este ciclo hasta la fecha de la prueba. Cada sesión dura aproximadamente
\MinutosSesion{} minutos; la meta es completar al menos \MetaSemanal{} sesiones por semana. Si un
día tienes menos tiempo, haz sólo la recuperación de memoria y una consulta corta.

\begin{longtable}{p{0.10\textwidth} p{0.25\textwidth} p{0.55\textwidth}}
\toprule
\textbf{Día} & \textbf{Foco} & \textbf{Producto comprobable} \\
\midrule
\endhead
1 & Método y esquema & Dibuja las relaciones entre empleados, cargos, departamentos y jefes sin
mirar la guía. \\
2 & Módulos 1--2 & Resuelve una consulta con concatenación y otra con filtros, nulos y orden. \\
3 & Módulos 3--4 & Resuelve una consulta con funciones de texto y otra con conversiones o nulos. \\
4 & Módulo 5 & Resuelve un resumen con \campo{GROUP BY} y justifica por qué el filtro usa
\campo{HAVING}. \\
5 & Módulo 6 & Construye un \campo{JOIN} y una auto-unión siguiendo las claves del esquema. \\
6 & Módulo 7 & Resuelve una subconsulta de una fila y otra de varias filas. \\
7 & Pauta completa & Simula los seis requerimientos con un límite de 90 minutos y corrige con la
rúbrica. \\
\bottomrule
\end{longtable}

\begin{rutina}
En cada sesión usa cuatro bloques: \textbf{5 minutos} para recordar sintaxis sin mirar,
\textbf{10 minutos} de lectura, \textbf{20 minutos} escribiendo SQL y \textbf{5 minutos} para
corregir. Marca cada error como conceptual, de sintaxis o de lectura del enunciado. Al día siguiente
comienza rehaciendo la consulta que fallaste.
\end{rutina}

\subsection{Registro de avance}

Marca una casilla sólo si puedes resolver el tema sin copiar el ejemplo. Usa V (verde) si lo haces
solo, A (amarillo) si necesitas una pista y R (rojo) si necesitas mirar la respuesta.

\begin{center}
\begin{tabularx}{\textwidth}{>{\bfseries}p{0.33\textwidth} *{3}{>{\centering\arraybackslash}X}}
\toprule
\textbf{Habilidad} & \textbf{Intento 1} & \textbf{Intento 2} & \textbf{Intento 3} \\
\midrule
SELECT, alias y concatenación & & & \\
WHERE, LIKE, rangos y nulos & & & \\
Funciones de texto, número y fecha & & & \\
GROUP BY y HAVING & & & \\
JOIN y auto-unión & & & \\
Subconsultas & & & \\
Pauta completa en 90 minutos & & & \\
\bottomrule
\end{tabularx}
\end{center}

\subsection{El método de seis decisiones}

Antes de escribir SQL, contesta en papel:

\begin{enumerate}[label=\arabic*.]
  \item \textbf{¿Qué columnas debe mostrar el resultado?} Eso construye el \campo{SELECT}.
  \item \textbf{¿De qué tablas salen?} Eso construye el \campo{FROM}.
  \item \textbf{¿Cómo se relacionan las tablas?} Eso construye cada \campo{JOIN ... ON}.
  \item \textbf{¿Qué filas individuales se aceptan?} Eso construye el \campo{WHERE}.
  \item \textbf{¿Hay que resumir filas?} Si hay conteos, promedios, mínimos, máximos o sumas,
  piensa en funciones de grupo y \campo{GROUP BY}.
  \item \textbf{¿Se filtran grupos o se ordena el resultado?} Usa \campo{HAVING} para grupos y
  \campo{ORDER BY} para presentación.
\end{enumerate}

El orden de escritura que usaremos durante toda la guía es:

\begin{lstlisting}
SELECT   expresiones
FROM     tablas_y_joins
WHERE    condiciones_de_fila
GROUP BY columnas_de_agrupacion
HAVING   condiciones_de_grupo
ORDER BY criterio_de_orden;
\end{lstlisting}

No todas las consultas necesitan todas las cláusulas, pero las que aparezcan deben conservar ese
orden.

\section{Plantillas textuales para construir consultas}

Esta sección no entrega consultas para copiar. Cada plantilla indica, en palabras, qué debes colocar
después de un comando. Los textos en cursiva representan decisiones que cambian según el enunciado;
no se escriben literalmente. Lee primero el requerimiento, reemplaza mentalmente cada decisión y
recién entonces redacta la consulta.

\begin{clave}
\campo{GROUP BY} no sirve para enlazar tablas. Su función es reunir filas que comparten el mismo
valor. Para enlazar una tabla con otra se utiliza \campo{JOIN} y la relación se explica con
\campo{ON} o, cuando corresponde, con \campo{USING}.
\end{clave}

\subsection{Consulta básica y construcción de la salida}

\begin{plantilla}{SELECT y FROM}
\textbf{Qué debes colocar.} Después de \campo{SELECT}, escribe los elementos que deben aparecer como
columnas del resultado. Cada elemento puede ser un atributo, un cálculo, una función o un texto fijo;
separa varios elementos con comas. Después de \campo{FROM}, escribe la tabla que contiene las filas
de partida y, si será útil, asígnale un alias corto.

\textbf{Frase para completar.} ``Mostrar \emph{los datos solicitados} desde
\emph{la tabla que contiene las filas principales}.''

\textbf{Paso a paso.}
\begin{enumerate}[label=\arabic*.]
  \item Subraya en el enunciado todo lo que se debe mostrar.
  \item Convierte cada dato solicitado en un elemento del \campo{SELECT}.
  \item Comprueba que cada elemento esté separado por una coma.
  \item Identifica la tabla que contiene las filas principales y colócala en \campo{FROM}.
  \item Si existen varias tablas, deja aquí la principal y añade las demás mediante \campo{JOIN}.
\end{enumerate}
\end{plantilla}

\begin{plantilla}{alias con AS}
Un alias cambia únicamente el encabezado mostrado o el nombre corto utilizado dentro de la consulta.
Después de una columna o expresión, coloca \campo{AS} y luego un nombre descriptivo. Si el alias
contiene espacios, tildes o debe conservar exactamente mayúsculas y minúsculas, escríbelo entre
comillas dobles. Para un alias de tabla se coloca el nombre corto inmediatamente después del nombre
de la tabla; en Oracle no es necesario escribir \campo{AS} para ese alias.
\end{plantilla}

\begin{plantilla}{DISTINCT}
Coloca \campo{DISTINCT} inmediatamente después de \campo{SELECT} cuando el resultado deba eliminar
filas repetidas. La comparación considera en conjunto todos los elementos mostrados. Antes de usarlo,
comprueba que la repetición no provenga de un \campo{JOIN} incorrecto: \campo{DISTINCT} limpia la
salida, pero no repara una relación mal construida.
\end{plantilla}

\begin{plantilla}{expresiones y concatenación}
Para un cálculo, coloca en el \campo{SELECT} el atributo numérico, el operador aritmético y el valor u
otro atributo con el que se operará. Usa paréntesis cuando el orden matemático deba quedar explícito.
Para formar una frase o un dato compuesto, intercala los atributos con los textos fijos y únelos con
el operador de concatenación. Los textos fijos llevan comillas simples; los nombres de columnas no.
Si utilizas la función \campo{CONCAT}, recuerda que recibe dos elementos por vez, por lo que más de
dos elementos requieren anidación.
\end{plantilla}

\subsection{Filtros de filas y condiciones}

\begin{plantilla}{WHERE y operadores de comparación}
Después de \campo{WHERE}, coloca primero el atributo que se evaluará, luego el operador de
comparación y finalmente el valor, atributo o resultado con el que se compara. Usa igualdad para una
coincidencia exacta; los operadores de diferencia, mayor, menor, mayor o igual y menor o igual se
eligen de acuerdo con las palabras del enunciado. Los textos y fechas convertidas explícitamente
requieren el tratamiento de su tipo; los números no llevan comillas.

\textbf{Frase para completar.} ``Conservar las filas donde \emph{el atributo evaluado} cumpla
\emph{la condición indicada}.''

\textbf{Paso a paso.}
\begin{enumerate}[label=\arabic*.]
  \item Identifica qué filas deben quedar dentro del resultado.
  \item Busca el atributo que permite decidir si cada fila se acepta.
  \item Traduce expresiones como ``igual'', ``distinto'', ``superior'' o ``hasta'' al operador
  correspondiente.
  \item Escribe el valor de comparación con el tipo correcto.
  \item Si hay más condiciones, conéctalas con \campo{AND} u \campo{OR} y usa paréntesis cuando pueda
  existir ambigüedad.
\end{enumerate}
\end{plantilla}

\begin{plantilla}{AND, OR y NOT}
Usa \campo{AND} cuando todas las condiciones conectadas deban cumplirse. Usa \campo{OR} cuando sea
suficiente cumplir al menos una. Coloca \campo{NOT} delante de la condición que se desea negar.
Como \campo{AND} se evalúa antes que \campo{OR}, agrupa con paréntesis las condiciones que expresan
una misma idea. Después de escribirlas, lee la condición completa en voz alta y verifica que diga lo
mismo que el requerimiento.
\end{plantilla}

\begin{plantilla}{BETWEEN y NOT BETWEEN}
Después de \campo{BETWEEN}, coloca el límite inferior, luego \campo{AND} y finalmente el límite
superior. Ambos extremos están incluidos. Usa \campo{NOT BETWEEN} cuando se pidan filas fuera del
rango; en ese caso tampoco se incluyen los dos extremos. Verifica que ambos límites tengan un tipo
compatible con el atributo evaluado.
\end{plantilla}

\begin{plantilla}{IN y NOT IN}
Después de \campo{IN}, coloca entre paréntesis el conjunto de valores aceptados o una subconsulta que
produzca esos valores. Es la opción natural cuando un atributo puede coincidir con cualquiera de
varias alternativas. \campo{NOT IN} expresa que no debe coincidir con ninguna, pero requiere cuidado
si el conjunto puede contener nulos; en ese caso conviene excluirlos dentro de la subconsulta o
evaluar otra estrategia.
\end{plantilla}

\begin{plantilla}{LIKE}
Después de \campo{LIKE}, coloca entre comillas simples el patrón de búsqueda. El símbolo de porcentaje
representa cero o más caracteres y el guion bajo representa exactamente un carácter. Coloca el
porcentaje al final para ``comienza con'', al inicio para ``termina con'' y a ambos lados para
``contiene''. Si la comparación no debe depender de mayúsculas, transforma tanto el atributo como el
patrón siguiendo el mismo criterio.
\end{plantilla}

\begin{plantilla}{IS NULL e IS NOT NULL}
Después del atributo coloca \campo{IS NULL} para conservar filas donde falta el dato, o
\campo{IS NOT NULL} para conservar aquellas donde sí existe. Un nulo no es cero, una cadena vacía ni
un valor que pueda compararse mediante igualdad. Primero interpreta si el enunciado pide ausencia o
presencia y luego elige una de estas dos condiciones.
\end{plantilla}

\subsection{Orden, funciones y transformación de valores}

\begin{plantilla}{ORDER BY}
Después de \campo{ORDER BY}, coloca el atributo, expresión o alias que determina el orden. Añade
\campo{ASC} para orden ascendente o \campo{DESC} para descendente; si no indicas nada, se usa el orden
ascendente. Para resolver empates, agrega más criterios separados por comas, desde el criterio más
importante hasta el último desempate. Esta cláusula organiza la presentación y se escribe al final.

\textbf{Frase para completar.} ``Ordenar primero por \emph{el criterio principal} en sentido
\emph{ascendente o descendente}, y luego por \emph{el criterio de desempate}.''
\end{plantilla}

\begin{plantilla}{funciones de una sola fila}
Coloca el nombre de la función y, entre paréntesis, el atributo o valor que debe transformar. Una
función de texto recibe texto y puede cambiar mayúsculas, extraer caracteres, medir longitud o buscar
una posición. Una función numérica recibe números y puede redondear, truncar u obtener un resto. Una
función de fecha recibe una fecha y permite desplazarla, encontrar límites del calendario o calcular
una diferencia. Estas funciones producen un resultado por cada fila y pueden anidarse: se resuelve
primero la función más interna. Cuando sólo necesites evaluar una expresión y no leer datos del
modelo, utiliza en \campo{FROM} la tabla auxiliar de una fila destinada a ese propósito.
\end{plantilla}

\begin{plantilla}{conversiones y formatos}
Para convertir explícitamente, elige la función según el tipo de destino: texto, número o fecha.
Dentro de sus paréntesis coloca primero el valor que se convertirá y después, cuando corresponda, el
modelo de formato. Calcula y compara usando los tipos originales; aplica el formato sólo cuando el
resultado esté listo para mostrarse. Esto evita ordenar números como palabras o interpretar fechas
según una configuración diferente.
\end{plantilla}

\begin{plantilla}{tratamiento de nulos}
Usa \campo{NVL} cuando necesites sustituir un posible nulo por un único valor alternativo. Usa
\campo{NVL2} cuando quieras un resultado si el dato existe y otro si falta. Usa \campo{COALESCE}
cuando debas elegir el primer valor no nulo de una lista. Usa \campo{NULLIF} cuando dos expresiones
iguales deban producir un nulo. Antes de aplicar cualquiera, decide qué significa la ausencia en el
problema y comprueba que los resultados sean de tipos compatibles.
\end{plantilla}

\begin{plantilla}{CASE y DECODE}
Para una decisión general, comienza una expresión \campo{CASE}, añade una condición con
\campo{WHEN}, coloca el resultado correspondiente después de \campo{THEN}, repite las ramas
necesarias, agrega un resultado alternativo con \campo{ELSE} y cierra con \campo{END}. Para
comparaciones simples de igualdad, \campo{DECODE} recibe una expresión y pares formados por valor de
búsqueda y resultado, más un resultado predeterminado opcional. Prefiere \campo{CASE} cuando existan
rangos, desigualdades o varias condiciones lógicas.
\end{plantilla}

\subsection{Resúmenes, GROUP BY y HAVING}

\begin{plantilla}{funciones de grupo}
Coloca una función de grupo alrededor del atributo que deseas resumir. Usa \campo{COUNT} para contar,
\campo{SUM} para sumar, \campo{AVG} para promediar, \campo{MIN} para obtener el menor valor y
\campo{MAX} para obtener el mayor. \campo{COUNT(*)} cuenta filas; una función aplicada a un atributo
normalmente ignora sus valores nulos. Si no existe \campo{GROUP BY}, la función resume todo el
conjunto filtrado en un único grupo.
\end{plantilla}

\begin{plantilla}{GROUP BY}
Después de \campo{GROUP BY}, coloca el atributo o los atributos cuyos valores iguales definen cada
grupo. Es decir, escribe aquello que responde ``por cada qué'' se desea calcular el resumen. Si el
resultado muestra varios atributos normales además de funciones de grupo, todos esos atributos
normales deben aparecer también en \campo{GROUP BY}. No coloques allí las funciones de grupo.

\textbf{Frase para completar.} ``Agrupar las filas por \emph{el atributo que identifica cada grupo}
para calcular \emph{el conteo, suma, promedio, mínimo o máximo solicitado}.''

\textbf{Paso a paso.}
\begin{enumerate}[label=\arabic*.]
  \item Encuentra en el enunciado expresiones como ``por cada'', ``por categoría'' o ``por tipo''.
  \item Coloca ese atributo de agrupación como dato visible si el resultado debe identificar cada
  grupo.
  \item Añade al resultado las funciones de resumen solicitadas.
  \item Aplica en \campo{WHERE} los filtros que eliminan filas antes de agrupar.
  \item Después de \campo{GROUP BY}, repite todos los atributos visibles que no estén dentro de una
  función de grupo.
  \item Comprueba que la salida produzca una fila por cada combinación distinta de esos atributos.
\end{enumerate}
\end{plantilla}

\begin{plantilla}{HAVING}
Después de \campo{HAVING}, coloca una condición que evalúe el resultado de una función de grupo o una
característica del grupo ya formado. Esta cláusula decide qué grupos permanecen. La secuencia mental
es: primero filtrar filas con \campo{WHERE}, después formar grupos con \campo{GROUP BY}, luego
calcular los resúmenes y finalmente filtrar esos grupos con \campo{HAVING}. Si la condición puede
evaluarse fila por fila antes de agrupar, normalmente corresponde a \campo{WHERE}.

\textbf{Frase para completar.} ``Conservar solamente los grupos cuyo
\emph{resultado agregado} cumpla \emph{la condición solicitada}.''
\end{plantilla}

\subsection{Uniones entre tablas}

\begin{plantilla}{JOIN con ON}
Después de la tabla principal, coloca el tipo de \campo{JOIN}, el nombre de la tabla secundaria y su
alias. A continuación escribe \campo{ON}; a su izquierda coloca el alias y la clave de una tabla, y a
su derecha el alias y la clave relacionada de la otra. Ambas columnas deben representar la misma
relación del modelo, no solamente compartir un nombre parecido.

\textbf{Frase para completar.} ``Unir \emph{la tabla auxiliar} con \emph{la tabla principal} cuando
\emph{la clave de la primera} sea igual a \emph{la clave relacionada de la segunda}.''

\textbf{Paso a paso.}
\begin{enumerate}[label=\arabic*.]
  \item Decide cuál es la entidad principal que determina las filas que quieres listar.
  \item Identifica qué dato falta y en qué tabla auxiliar se encuentra.
  \item Busca en el modelo la clave que sale de la tabla principal y la clave a la que apunta.
  \item Asigna un alias diferente a cada tabla.
  \item Escribe en \campo{ON} la igualdad entre ambas claves, calificadas con sus alias.
  \item Repite el proceso por cada tabla adicional y revisa si alguna unión multiplica filas de forma
  inesperada.
\end{enumerate}
\end{plantilla}

\begin{plantilla}{INNER, LEFT, RIGHT y FULL JOIN}
Usa \campo{INNER JOIN}, o simplemente \campo{JOIN}, cuando sólo deban permanecer filas con
coincidencia en ambas tablas. Usa \campo{LEFT JOIN} cuando deban conservarse todas las filas de la
tabla escrita a la izquierda, aunque no tengan coincidencia. \campo{RIGHT JOIN} conserva todas las
filas de la tabla derecha. \campo{FULL JOIN} conserva las filas de ambos lados. Primero identifica
qué conjunto debe preservarse y luego ubícalo en el lado coherente con el tipo de unión elegido.
\end{plantilla}

\begin{plantilla}{USING, NATURAL JOIN, auto-unión y cadena de uniones}
Usa \campo{USING} únicamente cuando las dos tablas comparten una columna de enlace con exactamente el
mismo nombre; dentro de sus paréntesis coloca ese nombre una sola vez. Para una auto-unión, escribe la
misma tabla dos veces con alias distintos, asignando a cada copia un rol diferente, y relaciona esos
roles mediante \campo{ON}. Para una cadena, añade una tabla a la vez siguiendo el camino real de
claves y verifica cada enlace antes de continuar con el siguiente. \campo{NATURAL JOIN} busca de
forma automática todas las columnas con nombres coincidentes; úsalo sólo si el modelo confirma que
todas esas coincidencias representan las relaciones deseadas, porque puede enlazar columnas
adicionales sin que resulte evidente.
\end{plantilla}

\begin{plantilla}{unión no equitativa y CROSS JOIN}
Una unión no equitativa utiliza en \campo{ON} una condición de rango o desigualdad en lugar de una
igualdad, cuando el modelo define la coincidencia de esa forma. \campo{CROSS JOIN} combina cada fila
de una tabla con todas las filas de la otra y no utiliza condición de enlace. Debe aparecer sólo
cuando el producto cartesiano sea solicitado expresamente; si surge por accidente, revisa las claves
de unión.
\end{plantilla}

\subsection{Subconsultas}

\begin{plantilla}{subconsulta de una fila}
Dentro de paréntesis, coloca una consulta interna que responda primero la pregunta auxiliar y devuelva
un único valor. En la consulta externa, compara un atributo con ese resultado mediante el operador
simple apropiado. Antes de usar igualdad, mayor o menor, asegúrate de que la consulta interna no pueda
devolver más de una fila; de lo contrario se producirá un error.

\textbf{Frase para completar.} ``Conservar las filas cuyo \emph{atributo externo} se compare con
\emph{el único valor calculado por la pregunta interna}.''
\end{plantilla}

\begin{plantilla}{subconsulta de varias filas}
Dentro de paréntesis, coloca una consulta interna que produzca una columna con cero, uno o varios
valores. En la consulta externa, utiliza un operador preparado para conjuntos, como \campo{IN}, para
comparar el atributo con cualquiera de esos resultados. La consulta interna debe mostrar sólo la
columna necesaria para la comparación y sus valores deben ser compatibles con el atributo externo.

\textbf{Frase para completar.} ``Conservar las filas cuyo \emph{atributo externo} pertenezca al
\emph{conjunto de valores obtenido por la pregunta interna}.''
\end{plantilla}

\begin{plantilla}{método para construir una subconsulta}
\begin{enumerate}[label=\arabic*.]
  \item Separa el requerimiento en una pregunta principal y una pregunta auxiliar.
  \item Resuelve primero la pregunta auxiliar y determina cuántas filas podría devolver.
  \item Coloca esa consulta entre paréntesis en la condición de la pregunta principal.
  \item Usa un operador simple para un solo valor o un operador de conjunto para varios valores.
  \item Controla los nulos cuando puedan alterar una comparación negativa.
  \item Ejecuta mentalmente primero la parte interna y después la externa para justificar el
  resultado.
\end{enumerate}
\end{plantilla}

\subsection{Plantilla maestra de lectura}

Para cualquier requerimiento, completa en una hoja las siguientes frases antes de escribir SQL:

\begin{enumerate}[label=\textbf{\arabic*.}]
  \item Mostraré \emph{los atributos, cálculos o textos pedidos}.
  \item La fila principal proviene de \emph{la tabla que representa lo que se desea listar}.
  \item Añadiré \emph{cada tabla auxiliar} mediante \emph{su clave de enlace real}.
  \item Conservaré únicamente las filas que cumplan \emph{las condiciones individuales}.
  \item Resumiré con \emph{la función solicitada} por cada \emph{atributo de agrupación}, si
  corresponde.
  \item Conservaré los grupos que cumplan \emph{la condición agregada}.
  \item Presentaré el resultado según \emph{el criterio y sentido de orden solicitado}.
\end{enumerate}

Cuando una frase no aplica, omite su cláusula. Las frases que sí aplican deben transformarse
respetando el orden de escritura de la sentencia.

\section{El esquema OEHR: aprender a seguir las claves}

Una consulta de varias tablas sólo es correcta si sus relaciones representan el modelo. Las tablas
más importantes para la prueba son las siguientes.

\subsection*{Cómo leer un esquema antes de consultar}

Una base de datos relacional separa la información en tablas para no repetir innecesariamente los
mismos datos. En OEHR, por ejemplo, la fila de un empleado no guarda el texto completo del
departamento ni el nombre de su cargo: guarda \campo{DEPARTMENT_ID} y \campo{JOB_ID}. Esos valores
funcionan como referencias hacia las filas correspondientes de \tabla{OEHR_DEPARTMENTS} y
\tabla{OEHR_JOBS}. Para resolver un requerimiento debes distinguir la \emph{tabla principal}, que
representa aquello que el enunciado quiere listar, de las tablas auxiliares que aportan atributos. Si
el enunciado dice ``muestre todos los trabajadores con su departamento'', la tabla principal es
\tabla{OEHR_EMPLOYEES}: quieres obtener una fila de salida por trabajador. Como el nombre del
departamento no está allí, sigues la referencia \campo{EMPLOYEES.DEPARTMENT_ID} hasta la clave
\campo{DEPARTMENTS.DEPARTMENT_ID}. La igualdad no se inventa por semejanza semántica, sino que se
deduce de la clave foránea declarada en el esquema. En esta relación, muchos empleados pueden
pertenecer a un mismo departamento, pero cada empleado apunta como máximo a un departamento. Por
eso, al unir correctamente ambas tablas, un empleado no debería multiplicarse por todos los
departamentos. Este razonamiento también permite anticipar nulos: \campo{DEPARTMENT_ID} admite nulo,
por lo que podría existir un empleado sin departamento asignado. Si el requerimiento exige conservarlo,
la unión tendrá que preservar la tabla principal mediante \campo{LEFT JOIN}; si sólo interesan las
coincidencias, bastará un \campo{JOIN} interno. Cuando la información solicitada está más lejos, no se
salta directamente entre tablas: se recorre el camino completo de claves. Para llegar desde empleado a
región se pasa por departamento, ubicación y país. Pensar el esquema como un mapa de nodos y enlaces
evita dos fallos comunes: unir columnas que no representan la misma entidad y ocultar un error con
\campo{DISTINCT}. \campo{DISTINCT} puede quitar duplicados visibles, pero jamás corrige una relación
mal construida.

\begin{longtable}{p{0.26\textwidth} p{0.65\textwidth}}
\toprule
\textbf{Tabla} & \textbf{Columnas relevantes} \\
\midrule
\endhead
\tabla{OEHR\_EMPLOYEES} &
\campo{EMPLOYEE\_ID}, \campo{FIRST\_NAME}, \campo{LAST\_NAME}, \campo{EMAIL},
\campo{HIRE\_DATE}, \campo{JOB\_ID}, \campo{SALARY}, \campo{COMMISSION\_PCT},
\campo{MANAGER\_ID}, \campo{DEPARTMENT\_ID}. \\
\tabla{OEHR\_JOBS} &
\campo{JOB\_ID}, \campo{JOB\_TITLE}, \campo{MIN\_SALARY}, \campo{MAX\_SALARY}. \\
\tabla{OEHR\_DEPARTMENTS} &
\campo{DEPARTMENT\_ID}, \campo{DEPARTMENT\_NAME}, \campo{MANAGER\_ID},
\campo{LOCATION\_ID}. \\
\tabla{OEHR\_LOCATIONS} &
\campo{LOCATION\_ID}, \campo{STREET\_ADDRESS}, \campo{POSTAL\_CODE}, \campo{CITY},
\campo{STATE\_PROVINCE}, \campo{COUNTRY\_ID}. \\
\tabla{OEHR\_COUNTRIES} &
\campo{COUNTRY\_ID}, \campo{COUNTRY\_NAME}, \campo{REGION\_ID}. \\
\tabla{OEHR\_REGIONS} &
\campo{REGION\_ID}, \campo{REGION\_NAME}. \\
\bottomrule
\end{longtable}

\subsection{Relaciones que debes reconocer}

\begin{align*}
\text{EMPLOYEES.JOB\_ID} &= \text{JOBS.JOB\_ID} \\
\text{EMPLOYEES.DEPARTMENT\_ID} &= \text{DEPARTMENTS.DEPARTMENT\_ID} \\
\text{EMPLOYEES.MANAGER\_ID} &= \text{EMPLOYEES.EMPLOYEE\_ID} \\
\text{DEPARTMENTS.LOCATION\_ID} &= \text{LOCATIONS.LOCATION\_ID} \\
\text{LOCATIONS.COUNTRY\_ID} &= \text{COUNTRIES.COUNTRY\_ID} \\
\text{COUNTRIES.REGION\_ID} &= \text{REGIONS.REGION\_ID}
\end{align*}

La tercera relación es especial: una fila de empleados apunta a otra fila de la misma tabla. Eso
permite encontrar al jefe mediante una \emph{auto-unión} o \emph{self join}.

\begin{pausa}
Sin escribir SQL, traza el camino necesario para mostrar el apellido de un empleado y el continente
donde trabaja. La ruta es: empleado $\rightarrow$ departamento $\rightarrow$ ubicación
$\rightarrow$ país $\rightarrow$ región.
\end{pausa}

\section{Módulo 1: SELECT, expresiones, alias y concatenación}

\begin{objetivo}
Construir la salida de una consulta de una tabla, calcular columnas y producir textos legibles.
\end{objetivo}

El primer paso de toda consulta consiste en definir con exactitud la forma del resultado. Una
sentencia \campo{SELECT} no modifica los datos: construye una tabla temporal de salida a partir de las
filas disponibles en \campo{FROM}. Cada elemento separado por coma dentro de \campo{SELECT} se
convierte en una columna del resultado y puede ser una columna almacenada, un texto fijo, una
operación aritmética o una función. Esta idea ayuda a traducir enunciados: si se solicita ``código,
nombre completo y sueldo anual'', aparecen tres expresiones, aunque el nombre completo use dos
columnas y el sueldo anual se calcule multiplicando por doce. El resultado no necesita copiar la
estructura física de la tabla. Puede presentar datos derivados y encabezados claros mediante alias.
Sin embargo, la consulta se evalúa para cada fila producida por \campo{FROM}; por eso una expresión
como \campo{salary*12} genera un valor anual para cada empleado, no un total de la empresa. Para
obtener totales se requieren funciones de grupo, estudiadas más adelante. También debes diferenciar
una constante de una columna: \campo{'Empleado: '} es texto literal porque está entre comillas
simples, mientras que \campo{last_name} es una referencia al dato de la fila actual. La concatenación
combina esas piezas para construir reportes humanos sin alterar los valores originales. Los
paréntesis son decisivos en cálculos: \campo{(salary+100)*12} aplica un bono mensual antes de
anualizar, mientras que \campo{salary+100*12} sólo suma 1200 al sueldo mensual por la precedencia de
la multiplicación. Finalmente, \campo{DISTINCT} debe entenderse sobre la fila de salida completa. Si
seleccionas únicamente \campo{job_id}, obtienes cargos distintos; si seleccionas además
\campo{employee_id}, cada empleado vuelve a diferenciar la fila. Antes de ejecutar, cuenta las
expresiones solicitadas, verifica las comas y pregúntate qué valor produciría cada una para un solo
empleado concreto. Esa simulación mental detecta muchos errores tempranos.

\subsection{SELECT y FROM}

\campo{SELECT} indica qué se mostrará y \campo{FROM} de dónde se obtendrán las filas.

\begin{lstlisting}
SELECT employee_id, first_name, salary
FROM oehr_employees;
\end{lstlisting}

El asterisco devuelve todas las columnas. Es útil para explorar, pero en una respuesta de prueba es
mejor nombrar únicamente las columnas solicitadas.

\begin{lstlisting}
SELECT *
FROM oehr_departments;
\end{lstlisting}

\subsection{Expresiones aritméticas y precedencia}

Las columnas numéricas pueden participar en expresiones con \campo{+}, \campo{-}, \campo{*} y
\campo{/}. Multiplicación y división se evalúan antes que suma y resta; usa paréntesis para expresar
la intención.

\begin{lstlisting}
SELECT employee_id,
       salary,
       salary * 12 AS sueldo_anual,
       (salary + 57) * 12 AS anual_con_bono_mensual
FROM oehr_employees;
\end{lstlisting}

Un nulo propagará normalmente el nulo en una expresión. Por ejemplo,
\campo{salary * commission\_pct} será nulo cuando la comisión sea nula. Más adelante se resolverá
con \campo{NVL}.

\subsection{Alias de columna}

Un alias cambia el encabezado del resultado, no el nombre almacenado de la columna. Puede escribirse
con \campo{AS}. Las comillas dobles conservan espacios, tildes o mayúsculas exactas.

\begin{lstlisting}
SELECT employee_id AS "COD_EMPLEADO",
       salary * 12 AS "SUELDO_ANUAL"
FROM oehr_employees;
\end{lstlisting}

No uses comillas simples para un alias: en Oracle las comillas simples representan texto literal.

\subsection{DISTINCT}

\campo{DISTINCT} elimina filas duplicadas de la combinación completa de expresiones seleccionadas.

\begin{lstlisting}
SELECT DISTINCT job_id
FROM oehr_employees;
\end{lstlisting}

\subsection{Concatenación}

Oracle permite concatenar con \campo{||}. También ofrece \campo{CONCAT}, pero ésta recibe dos
argumentos; para unir tres piezas debe anidarse.

\begin{lstlisting}
SELECT employee_id,
       first_name || ' ' || last_name AS nombre_completo
FROM oehr_employees;

SELECT employee_id,
       CONCAT(first_name, CONCAT(' ', last_name)) AS nombre_completo
FROM oehr_employees;
\end{lstlisting}

\begin{prueba}
La primera consulta de la pauta pide formar una frase con \campo{JOB\_TITLE} y \campo{JOB\_ID}.
La construcción más clara es:
\begin{lstlisting}
SELECT 'El codigo del cargo ' || job_title ||
       ' es ' || job_id AS reporte
FROM oehr_jobs;
\end{lstlisting}
Después se añadirá el filtro que exige que el título o el código comience con S.
\end{prueba}

\subsection{Lectura activa del módulo 1}

\begin{enumerate}[label=\textbf{M1.\arabic*}]
  \item \nivel{Muy fácil.} ¿Qué cláusula indica las columnas que se mostrarán?
  \item \nivel{Muy fácil.} ¿Qué símbolo permite pedir todas las columnas?
  \item \nivel{Muy fácil.} ¿Qué palabra elimina resultados duplicados?
  \item \nivel{Intermedia.} Escribe una consulta que muestre código, apellido y sueldo anual de cada
  empleado, con los alias \campo{CODIGO}, \campo{APELLIDO} y \campo{ANUAL}.
  \item \nivel{Intermedia.} Explica la diferencia entre \campo{salary + 100 * 12} y
  \campo{(salary + 100) * 12}.
  \item \nivel{Difícil.} Genera una sola columna con el texto
  ``KING gana 24000 al año'', usando el apellido real, el sueldo real y una expresión anual. No
  escribas valores fijos salvo los textos de unión.
\end{enumerate}

\section{Módulo 2: WHERE, operadores, patrones, nulos y orden}

\begin{objetivo}
Seleccionar filas individuales con precisión y presentar el resultado en el orden solicitado.
\end{objetivo}

Una consulta suele comenzar con muchas filas posibles y luego descartar las que no cumplen el
requerimiento. \campo{WHERE} expresa precisamente esa decisión fila por fila: para cada empleado,
Oracle evalúa una condición cuyo resultado lógico puede ser verdadero, falso o desconocido. Sólo las
filas verdaderas continúan. La tercera posibilidad aparece principalmente con nulos y explica por qué
\campo{commission_pct = NULL} no funciona: un valor desconocido no puede afirmarse igual a otro
valor desconocido; debe preguntarse \campo{IS NULL}. Al combinar condiciones conviene traducir el
español de manera literal. ``Empleados del departamento 50 u 80, y con sueldo mayor a 5000'' se
convierte en \campo{department_id IN (50,80) AND salary>5000}. Los paréntesis protegen la intención
cuando existe una mezcla de \campo{AND} y \campo{OR}; sin ellos, la precedencia puede aceptar filas
que el enunciado no quería. Los operadores abreviados también tienen significado exacto:
\campo{BETWEEN} incluye los dos límites e \campo{IN} representa una lista de alternativas. En texto,
\campo{LIKE} compara un patrón, no una frase aproximada. El porcentaje acepta cualquier cantidad de
caracteres y el guion bajo exactamente uno, de modo que \campo{'\_a\%s'} exige una \campo{a} en la
segunda posición y una \campo{s} al final. Para no depender de cómo estén escritas las mayúsculas,
transforma tanto el razonamiento como el patrón a un mismo caso, por ejemplo
\campo{UPPER(last_name) LIKE 'S\%'}. Una vez elegidas las filas, \campo{ORDER BY} modifica solamente
su presentación: no cambia qué filas existen ni arregla un filtro equivocado. Cuando se ordena por
varias columnas, la primera define los grupos principales de orden y las siguientes resuelven los
empates. Leer una consulta correctamente significa imaginar esta secuencia: obtener filas, evaluar
cada condición, conservar las verdaderas y recién entonces ordenar lo que quedó.

\subsection{Comparaciones y texto literal}

\campo{WHERE} conserva las filas cuya condición resulta verdadera. Los textos y fechas escritas como
literales se rodean con comillas simples; los números no.

\begin{lstlisting}
SELECT employee_id, last_name, salary
FROM oehr_employees
WHERE salary >= 10000;

SELECT employee_id, last_name
FROM oehr_employees
WHERE last_name = 'Whalen';
\end{lstlisting}

Operadores comunes: \campo{=}, \campo{<>}, \campo{>}, \campo{<}, \campo{>=} y \campo{<=}.

\subsection{AND, OR, NOT y precedencia}

\campo{AND} exige que ambas condiciones sean verdaderas; \campo{OR} exige al menos una;
\campo{NOT} niega. Oracle evalúa \campo{NOT}, luego \campo{AND} y luego \campo{OR}. Usa paréntesis
cuando el requerimiento mezcle operadores.

\begin{lstlisting}
SELECT employee_id, last_name, job_id, salary
FROM oehr_employees
WHERE (job_id = 'AD_PRES' OR job_id = 'SA_REP')
  AND salary > 15000;
\end{lstlisting}

\subsection{BETWEEN e IN}

\campo{BETWEEN} incluye ambos extremos. \campo{IN} compara con un conjunto de alternativas.

\begin{lstlisting}
SELECT employee_id, last_name, salary
FROM oehr_employees
WHERE salary BETWEEN 3000 AND 5000;

SELECT employee_id, last_name, department_id
FROM oehr_employees
WHERE department_id IN (50, 80, 90);
\end{lstlisting}

Para excluir un conjunto usa \campo{NOT IN}, no una cadena de desigualdades unida con \campo{OR}.
La expresión \campo{x <> 100 OR x <> 120} es prácticamente siempre verdadera.

\subsection{LIKE y sus comodines}

\campo{LIKE} compara patrones de texto:

\begin{itemize}
  \item \campo{\%} representa cero o más caracteres;
  \item \campo{\_} representa exactamente un carácter.
\end{itemize}

\begin{lstlisting}
-- Apellido que comienza con S
WHERE last_name LIKE 'S%'

-- Nombre de departamento cuya segunda letra es a y termina en s
WHERE LOWER(department_name) LIKE '_a%s'

-- Cargo cuya penultima letra es a
WHERE LOWER(job_id) LIKE '%a_'
\end{lstlisting}

Las comparaciones de texto distinguen mayúsculas y minúsculas según los valores y la configuración.
Una técnica del curso es normalizar el dato con \campo{UPPER} o \campo{LOWER} y comparar con un patrón
en el mismo caso.

\begin{lstlisting}
WHERE UPPER(job_title) LIKE 'S%'
   OR UPPER(job_id)    LIKE 'S%'
\end{lstlisting}

\subsection{NULL no es cero ni texto vacío}

Un nulo representa un valor desconocido, ausente o no aplicable. No se evalúa con
\campo{= NULL} ni \campo{<> NULL}.

\begin{lstlisting}
SELECT employee_id, last_name
FROM oehr_employees
WHERE commission_pct IS NULL;
\end{lstlisting}

Usa \campo{IS NOT NULL} para exigir un valor conocido.

\subsection{ORDER BY}

\campo{ORDER BY} se escribe al final. El orden predeterminado es ascendente; \campo{DESC} lo invierte.
Se puede ordenar por columna, alias o posición, aunque el nombre o alias suele ser más legible.

\begin{lstlisting}
SELECT employee_id, last_name, department_id, hire_date
FROM oehr_employees
WHERE department_id IN (50, 80)
ORDER BY department_id ASC, hire_date DESC, last_name ASC;
\end{lstlisting}

\begin{errorbox}
No coloques punto y coma antes de \campo{ORDER BY}. El punto y coma termina la sentencia completa.
\end{errorbox}

\subsection{Lectura activa del módulo 2}

\begin{enumerate}[label=\textbf{M2.\arabic*}]
  \item \nivel{Muy fácil.} ¿Qué operador comprueba si una columna es nula?
  \item \nivel{Muy fácil.} En \campo{LIKE}, ¿qué representa \campo{\_}?
  \item \nivel{Muy fácil.} ¿Qué palabra ordena de mayor a menor?
  \item \nivel{Intermedia.} Muestra empleados de los departamentos 50 u 80 cuyo sueldo esté entre
  4000 y 9000, incluidos los límites.
  \item \nivel{Intermedia.} Escribe el patrón para apellidos cuya segunda letra sea \campo{a} y cuya
  última letra sea \campo{s}, ignorando mayúsculas.
  \item \nivel{Difícil.} Muestra empleados que pertenezcan a los departamentos 50, 80 o 90 y que
  además no tengan comisión, ordenados por departamento ascendente y sueldo descendente.
\end{enumerate}

\section{Módulo 3: funciones de una sola fila}

\begin{objetivo}
Transformar texto, números y fechas sin cambiar la cantidad de filas de entrada.
\end{objetivo}

Las funciones de una sola fila son pequeñas transformaciones aplicadas independientemente a cada
registro. Si entran 107 empleados y no existe un filtro, una expresión con \campo{UPPER},
\campo{SUBSTR} o \campo{ROUND} sigue produciendo 107 resultados; no resume el conjunto. Para decidir
qué función usar, identifica primero el tipo de dato que tienes y el tipo de resultado que necesitas.
Con texto puedes cambiar capitalización, medir longitud, buscar una posición o extraer un fragmento.
Por ejemplo, crear un correo exige descomponer el requerimiento en piezas: extraer la primera letra
con \campo{SUBSTR(first_name,1,1)}, añadir el apellido y finalmente el dominio. Con números,
\campo{ROUND} y \campo{TRUNC} responden a preguntas diferentes: el primero aproxima al valor más
cercano y el segundo elimina las posiciones que sobran. Con fechas, Oracle almacena un valor de fecha,
no simplemente el texto que ves en pantalla. Por ello puede restar fechas para obtener días, sumar
días o emplear funciones que navegan por meses. Si calculas semanas trabajadas, la resta
\campo{SYSDATE-hire_date} entrega días, la división por siete los transforma en semanas y
\campo{TRUNC} elimina la parte de semana incompleta. Este razonamiento por unidades evita usar
funciones al azar. \tabla{DUAL} aparece cuando deseas evaluar una expresión sin recorrer una tabla de
negocio: permite comprobar cuánto devuelve \campo{ROUND(45.5674,2)} o cuál es la fecha actual. Las
funciones pueden anidarse, y Oracle evalúa desde la más interna hacia afuera. En
\campo{UPPER(SUBSTR(first_name,1,1))}, primero se extrae y luego se convierte. Antes de escribir una
expresión compleja, dibuja sus pasos y comprueba el tipo que sale de cada uno. Una función correcta
con argumentos incorrectos puede producir un resultado válido sintácticamente, pero distinto del que
el problema solicita.

Una función de una sola fila recibe valores de una fila y devuelve un resultado por esa fila. Puede
usarse en \campo{SELECT}, \campo{WHERE} y \campo{ORDER BY}, respetando los tipos de datos.

\subsection{Funciones de caracteres}

\begin{longtable}{p{0.25\textwidth} p{0.65\textwidth}}
\toprule
\textbf{Función} & \textbf{Idea y ejemplo} \\
\midrule
\endhead
\campo{UPPER(texto)} & Convierte a mayúsculas: \campo{UPPER(last\_name)}. \\
\campo{LOWER(texto)} & Convierte a minúsculas: \campo{LOWER(email)}. \\
\campo{INITCAP(texto)} & Coloca en mayúscula la inicial de cada palabra. \\
\campo{CONCAT(a,b)} & Une dos valores; para más valores, anida o usa \campo{||}. \\
\campo{SUBSTR(t,p,n)} & Extrae \campo{n} caracteres de \campo{t} desde la posición \campo{p}. Una
posición negativa cuenta desde el final. \\
\campo{LENGTH(texto)} & Devuelve la cantidad de caracteres. \\
\campo{INSTR(t,b)} & Devuelve la posición donde comienza \campo{b} dentro de \campo{t}; si no lo
encuentra, devuelve cero. \\
\bottomrule
\end{longtable}

Ejemplo del estilo de la clase: crear un correo con la primera letra del nombre, el apellido y un
dominio.

\begin{lstlisting}
SELECT employee_id,
       UPPER(SUBSTR(first_name, 1, 1) ||
             last_name || '@SQL.CL') AS email_propuesto
FROM oehr_employees;
\end{lstlisting}

Para extraer el último carácter se puede usar \campo{SUBSTR(texto,-1,1)}. Para quitar el penúltimo
carácter conviene razonar en segmentos: parte anterior, último carácter y dominio.

\subsection{Funciones numéricas y DUAL}

\campo{ROUND} redondea, \campo{TRUNC} corta sin redondear y \campo{MOD} calcula el resto de una
división.

\begin{lstlisting}
SELECT ROUND(45.5674, 2) AS redondeado,
       TRUNC(45.5674, 2) AS truncado,
       MOD(54, 5)        AS resto
FROM dual;
\end{lstlisting}

Resultado conceptual: \campo{45.57}, \campo{45.56} y \campo{4}. \tabla{DUAL} es una tabla especial
de una fila útil cuando la expresión no depende de datos de una tabla del usuario.

\begin{errorbox}
La sintaxis correcta es \campo{MOD(54,3)}, no \campo{MOD(54/3)}. En una sesión configurada con punto
decimal, escribe \campo{45.5674}; una coma separaría argumentos.
\end{errorbox}

\subsection{Aritmética y funciones de fecha}

En Oracle, sumar un número a una fecha agrega días y restar dos fechas produce una diferencia en
días. \campo{SYSDATE} devuelve la fecha y hora actuales del servidor.

\begin{lstlisting}
SELECT employee_id,
       hire_date,
       TRUNC((SYSDATE - hire_date) / 7) AS semanas_trabajadas
FROM oehr_employees
WHERE department_id IN (50, 80)
ORDER BY semanas_trabajadas;
\end{lstlisting}

Otras funciones del material:

\begin{itemize}
  \item \campo{MONTHS\_BETWEEN(fecha1,fecha2)}: diferencia aproximada en meses;
  \item \campo{ADD\_MONTHS(fecha,n)}: suma o resta meses;
  \item \campo{NEXT\_DAY(fecha,'MONDAY')}: siguiente día indicado posterior a la fecha;
  \item \campo{LAST\_DAY(fecha)}: último día del mes;
  \item \campo{ROUND(fecha,formato)} y \campo{TRUNC(fecha,formato)}: redondean o truncan fechas.
\end{itemize}

\begin{lstlisting}
SELECT employee_id,
       hire_date,
       ADD_MONTHS(hire_date, 6) AS revision,
       LAST_DAY(hire_date)      AS fin_mes_contrato
FROM oehr_employees;
\end{lstlisting}

Los nombres de días empleados por \campo{NEXT\_DAY} pueden depender del idioma de la sesión. En una
prueba, sigue la configuración indicada por el profesor.

\subsection{Lectura activa del módulo 3}

\begin{enumerate}[label=\textbf{M3.\arabic*}]
  \item \nivel{Muy fácil.} ¿Qué función convierte un texto a mayúsculas?
  \item \nivel{Muy fácil.} ¿Qué diferencia esencial existe entre \campo{ROUND} y \campo{TRUNC}?
  \item \nivel{Muy fácil.} ¿Para qué se usa \tabla{DUAL}?
  \item \nivel{Intermedia.} Construye un correo formado por las dos primeras letras del nombre, un
  punto, el apellido y \campo{@ucm.cl}, todo en minúsculas.
  \item \nivel{Intermedia.} Calcula los años completos aproximados trabajados dividiendo los meses
  obtenidos con \campo{MONTHS\_BETWEEN} por 12 y truncando el resultado.
  \item \nivel{Difícil.} Genera el identificador formado por la última letra del apellido, la primera
  letra del nombre y la longitud total de ambos nombres concatenados, todo en mayúsculas.
\end{enumerate}

\section{Módulo 4: conversiones, formatos, nulos y decisiones}

\begin{objetivo}
Controlar explícitamente los tipos de datos, mostrar fechas y dinero con formato y producir valores
alternativos cuando existan nulos.
\end{objetivo}

Los datos poseen tipos y Oracle necesita saber qué operaciones tienen sentido para cada tipo. Un
número puede sumarse, una fecha puede compararse cronológicamente y un texto puede concatenarse. A
veces Oracle intenta convertir automáticamente un tipo en otro, pero esa conversión implícita depende
de la configuración de la sesión. El texto \campo{'01/02/2025'} podría interpretarse como primero de
febrero o dos de enero; por eso una consulta estable usa \campo{TO_DATE} con un formato explícito.
\campo{TO_CHAR} realiza el camino contrario cuando el objetivo es presentar el valor, no seguir
calculando con él. Formatear \campo{salary} como dinero produce texto visual; si después necesitas
promediar, debes conservar el número original para el cálculo y formatear sólo el resultado final. Los
nulos exigen otra decisión. \campo{NVL(commission_pct,0)} afirma una regla de negocio para esa
expresión: cuando la comisión falta, será tratada como cero. Eso permite multiplicar el sueldo sin que
el nulo se propague. No significa que la base de datos haya guardado cero ni que ambos conceptos sean
siempre equivalentes. \campo{NVL2} y \campo{CASE} permiten producir mensajes diferentes según la
existencia o el valor de un dato. En una rama condicional, todos los resultados deben poder convertirse
a tipos compatibles. Si una rama devuelve dinero formateado y otra el texto \campo{'SIN COMISION'},
conviene que ambas devuelvan texto: se aplica \campo{TO_CHAR} al monto en la primera rama. La regla
práctica es separar tres fases: calcular usando números o fechas reales, decidir usando condiciones y
formatear al final para mostrar. Esta separación evita comparar fechas como palabras, ordenar montos
por su símbolo monetario o sumar cadenas que sólo parecen números.

\subsection{Conversión explícita}

Oracle puede hacer algunas conversiones implícitas, pero dependen de la configuración y pueden
fallar. Una solución robusta expresa la conversión.

\begin{longtable}{p{0.31\textwidth} p{0.59\textwidth}}
\toprule
\textbf{Función} & \textbf{Uso} \\
\midrule
\endhead
\campo{TO\_CHAR(fecha,formato)} & Convierte una fecha en texto controlando su representación. \\
\campo{TO\_CHAR(número,formato)} & Convierte un número en texto, por ejemplo para moneda. \\
\campo{TO\_NUMBER(texto,formato)} & Convierte texto numérico en número. \\
\campo{TO\_DATE(texto,formato)} & Interpreta texto como fecha usando un modelo explícito. \\
\bottomrule
\end{longtable}

\begin{lstlisting}
SELECT employee_id,
       TO_CHAR(hire_date, 'DD "DE" MONTH "DE" YYYY') AS fecha,
       TO_CHAR(salary, 'FM$999,999,990.00')           AS sueldo
FROM oehr_employees;
\end{lstlisting}

\campo{FM} reduce los espacios de relleno del formato. Las palabras literales dentro de un formato de
fecha se rodean con comillas dobles dentro del texto SQL.

\begin{lstlisting}
SELECT employee_id, last_name, hire_date
FROM oehr_employees
WHERE hire_date >= TO_DATE('2016-01-01', 'YYYY-MM-DD')
ORDER BY hire_date;
\end{lstlisting}

Esto es preferible a comparar con \campo{'01/01/2016'}, porque el significado de ese texto depende de
\campo{NLS\_DATE\_FORMAT}.

\subsection{NVL y NVL2}

\campo{NVL(valor,reemplazo)} devuelve el reemplazo sólo si el valor es nulo. Ambos deben ser de tipos
compatibles.

\begin{lstlisting}
SELECT employee_id,
       salary,
       NVL(commission_pct, 0) AS comision,
       salary * NVL(commission_pct, 0) AS monto_comision
FROM oehr_employees;
\end{lstlisting}

\campo{NVL2(valor,si_no_es_nulo,si_es_nulo)} elige entre dos resultados.

\begin{lstlisting}
SELECT employee_id,
       NVL2(commission_pct,
            TO_CHAR(commission_pct * 100, 'FM990') || '%',
            'SIN COMISION') AS estado_comision
FROM oehr_employees;
\end{lstlisting}

\subsection{COALESCE y NULLIF}

\campo{COALESCE(a,b,c)} devuelve el primer valor no nulo. \campo{NULLIF(a,b)} devuelve nulo cuando
ambos valores son iguales; en caso contrario devuelve \campo{a}.

\begin{lstlisting}
SELECT employee_id,
       COALESCE(TO_CHAR(commission_pct),
                TO_CHAR(manager_id),
                'SIN COMISION NI JEFE') AS dato_disponible
FROM oehr_employees;
\end{lstlisting}

\subsection{CASE y DECODE}

\campo{CASE} permite expresar decisiones legibles y es parte del estándar SQL.

\begin{lstlisting}
SELECT employee_id, salary,
       CASE
         WHEN salary >= 15000 THEN 'ALTO'
         WHEN salary >= 8000  THEN 'MEDIO'
         ELSE 'BAJO'
       END AS tramo_sueldo
FROM oehr_employees;
\end{lstlisting}

\campo{DECODE} es una función propia de Oracle útil para equivalencias simples.

\begin{lstlisting}
SELECT employee_id, department_id,
       DECODE(department_id,
              10, 'ADMINISTRACION',
              20, 'MARKETING',
              'OTRO') AS area
FROM oehr_employees;
\end{lstlisting}

Para rangos y condiciones compuestas suele ser más claro \campo{CASE}. Para una equivalencia corta,
\campo{DECODE} puede ser suficiente.

\subsection{Lectura activa del módulo 4}

\begin{enumerate}[label=\textbf{M4.\arabic*}]
  \item \nivel{Muy fácil.} ¿Qué función convierte una fecha en texto formateado?
  \item \nivel{Muy fácil.} ¿Qué devuelve \campo{NVL(NULL,0)}?
  \item \nivel{Muy fácil.} ¿Qué función devuelve el primer argumento no nulo de una lista?
  \item \nivel{Intermedia.} Muestra el sueldo con dos decimales y la fecha de contratación con el
  formato \campo{YYYY-MM-DD}.
  \item \nivel{Intermedia.} Muestra \campo{CON COMISION} o \campo{SIN COMISION} según
  \campo{commission\_pct}, utilizando \campo{CASE}.
  \item \nivel{Difícil.} Muestra nombre completo, sueldo con formato monetario y monto de comisión
  con formato monetario; quienes no tengan comisión deben mostrar \campo{SIN COMISION} en la última
  columna. Evita mezclar números y textos incompatibles.
\end{enumerate}

\section{Módulo 5: funciones de grupo, GROUP BY y HAVING}

\begin{objetivo}
Resumir muchas filas en resultados por empresa, departamento o cargo, y filtrar esos resultados.
\end{objetivo}

Agrupar significa cambiar la unidad de análisis. En una consulta común, la unidad es una fila de
empleado; después de \campo{GROUP BY department_id}, la unidad pasa a ser un departamento compuesto
por todas sus filas de empleados. Las funciones de grupo responden preguntas sobre ese conjunto:
cuántos miembros tiene, cuánto suman sus sueldos, cuál es el promedio o cuáles son los extremos. Esta
transformación explica la regla de \campo{GROUP BY}: si el resultado representa una fila por
departamento, Oracle no puede mostrar libremente \campo{last_name}, porque dentro del grupo existen
muchos apellidos y no hay una única respuesta. Debes agrupar también por apellido, resumirlo con una
función apropiada o retirarlo. El orden conceptual ayuda a ubicar filtros. \campo{WHERE} actúa cuando
todavía existen empleados individuales; por eso sirve para limitar el estudio a departamentos 50, 80
y 90 antes de agrupar. \campo{HAVING} actúa cuando ya existe cada grupo y sus valores calculados; por
eso puede preguntar si \campo{AVG(salary)>8000} o si \campo{COUNT(*)>=2}. En la pauta, la diferencia
\campo{MAX(salary)-MIN(salary)} sólo existe después de reunir los empleados de cada cargo, de modo que
debe escribirse en \campo{HAVING}. También hay que pensar cuidadosamente en nulos:
\campo{COUNT(*)} cuenta filas, mientras \campo{COUNT(commission_pct)} cuenta únicamente valores de
comisión conocidos. Del mismo modo, \campo{AVG(commission_pct)} ignora empleados sin comisión, pero
\campo{AVG(NVL(commission_pct,0))} los incorpora como cero y responde otra pregunta. Para verificar
una consulta agrupada, toma un grupo imaginario de tres filas, calcula manualmente la función y
pregunta si cada columna visible posee exactamente un valor para ese grupo. Si no lo posee, el diseño
del \campo{SELECT} y del \campo{GROUP BY} aún no está completo.

\subsection{Funciones de grupo}

\begin{itemize}
  \item \campo{COUNT(*)}: cuenta filas.
  \item \campo{COUNT(columna)}: cuenta valores no nulos de esa columna.
  \item \campo{COUNT(DISTINCT columna)}: cuenta valores diferentes no nulos.
  \item \campo{SUM(columna)}: suma valores no nulos.
  \item \campo{AVG(columna)}: promedio de valores no nulos.
  \item \campo{MIN(columna)} y \campo{MAX(columna)}: extremos del conjunto.
\end{itemize}

\begin{lstlisting}
SELECT ROUND(AVG(salary), 0) AS promedio,
       MIN(hire_date)        AS primera_contratacion,
       MAX(hire_date)        AS ultima_contratacion,
       COUNT(*)              AS cantidad_filas
FROM oehr_employees;
\end{lstlisting}

Las funciones ignoran nulos salvo \campo{COUNT(*)}, que cuenta filas. Esto importa al promediar
comisiones. \campo{AVG(commission\_pct)} promedia únicamente empleados con comisión;
\campo{AVG(NVL(commission\_pct,0))} incluye al resto como cero. No significan lo mismo.

\subsection{GROUP BY}

Sin \campo{GROUP BY}, la función resume todas las filas filtradas en un solo grupo. Con
\campo{GROUP BY}, produce un resultado por cada combinación diferente de columnas agrupadas.

\begin{lstlisting}
SELECT department_id,
       ROUND(AVG(salary), 0) AS promedio_sueldo
FROM oehr_employees
GROUP BY department_id
ORDER BY department_id;
\end{lstlisting}

Regla esencial: toda expresión del \campo{SELECT} que no sea agregada debe ser compatible con las
expresiones del \campo{GROUP BY}.

\begin{lstlisting}
SELECT department_id,
       job_id,
       TO_CHAR(SUM(salary), 'FM$999,999,990.00') AS total
FROM oehr_employees
GROUP BY department_id, job_id
ORDER BY department_id, job_id;
\end{lstlisting}

\subsection{WHERE frente a HAVING}

\campo{WHERE} filtra filas antes del agrupamiento; \campo{HAVING} filtra grupos después de calcular
las funciones de grupo.

\begin{lstlisting}
SELECT department_id,
       MAX(salary) AS sueldo_maximo
FROM oehr_employees
WHERE department_id BETWEEN 50 AND 90
GROUP BY department_id
HAVING MAX(salary) > 10000
ORDER BY department_id DESC;
\end{lstlisting}

Lectura paso a paso:

\begin{enumerate}
  \item se conservan empleados de departamentos 50 a 90;
  \item se forman grupos por departamento;
  \item se calcula el máximo de cada grupo;
  \item se conservan grupos cuyo máximo supera 10000;
  \item se ordena el resultado.
\end{enumerate}

\begin{prueba}
Una consulta de la pauta intenta colocar \campo{MAX(salary)-MIN(salary) <> 0} en \campo{WHERE}. Eso es
incorrecto porque la condición sólo existe después de formar el grupo. La forma conceptual correcta
es:
\begin{lstlisting}
SELECT job_id,
       COUNT(employee_id) AS cantidad_empleados,
       ROUND(AVG(salary), 2) AS promedio_sueldo
FROM oehr_employees
GROUP BY job_id
HAVING MAX(salary) - MIN(salary) <> 0;
\end{lstlisting}
\end{prueba}

\subsection{Lectura activa del módulo 5}

\begin{enumerate}[label=\textbf{M5.\arabic*}]
  \item \nivel{Muy fácil.} ¿Qué función calcula un promedio?
  \item \nivel{Muy fácil.} ¿Qué diferencia hay entre \campo{COUNT(*)} y
  \campo{COUNT(commission\_pct)}?
  \item \nivel{Muy fácil.} ¿Qué cláusula filtra grupos?
  \item \nivel{Intermedia.} Muestra cada \campo{job\_id} y la cantidad de empleados que lo poseen.
  \item \nivel{Intermedia.} Muestra los departamentos cuyo sueldo promedio es mayor que 8000.
  \item \nivel{Difícil.} Para los departamentos 50, 80 y 90, muestra departamento, cargo, cantidad
  de empleados y promedio de sueldo. Conserva sólo combinaciones con al menos dos empleados y ordena
  por promedio descendente.
\end{enumerate}

\section{Módulo 6: consultas de varias tablas y JOIN}

\begin{objetivo}
Combinar tablas sin producir resultados inventados, usando claves, alias y el tipo de unión correcto.
\end{objetivo}

Un \campo{JOIN} no consiste simplemente en colocar dos tablas juntas: define qué pares de filas
representan una misma relación del mundo modelado. Supón que el requerimiento pide código del
trabajador, nombre completo y nombre del departamento. Primero eliges
\tabla{OEHR_EMPLOYEES} como tabla principal porque el sujeto del listado es el trabajador. Allí están
el código y los nombres, pero sólo aparece \campo{DEPARTMENT_ID}; el texto
\campo{DEPARTMENT_NAME} vive en \tabla{OEHR_DEPARTMENTS}. Entonces asignas alias para expresar roles,
\campo{e} para empleado y \campo{d} para departamento, y conectas la referencia del empleado con la
clave del departamento: \campo{e.department_id=d.department_id}. Oracle toma una fila de empleado,
busca la fila departamental cuyo código coincide y forma una fila combinada que contiene columnas de
ambas. Después \campo{SELECT} puede escoger \campo{e.employee_id}, los nombres de \campo{e} y
\campo{d.department_name}. Si omites la condición, cada empleado se combina con cada departamento;
si escribes \campo{department_id=department_id} sin alias, la intención es ambigua; y si conectas
\campo{e.manager_id} con \campo{d.manager_id}, estarías respondiendo otra pregunta, no la pertenencia
del empleado a un departamento. El tipo de unión depende de qué filas deben sobrevivir. Un
\campo{INNER JOIN} conserva sólo empleados con departamento coincidente. Un
\campo{LEFT JOIN}, colocando empleados a la izquierda, conserva además al trabajador sin
departamento y completa las columnas de \campo{d} con nulos. Antes de elegir, busca expresiones como
``todos los empleados'', ``aunque no tengan'' o ``sólo los que poseen''. Finalmente, anticipa la
cardinalidad: muchos empleados pueden coincidir con un departamento, de modo que el departamento se
repite legítimamente en varias filas. Esa repetición no es un duplicado erróneo; representa distintos
empleados relacionados con la misma entidad.

\subsection{Por qué aparece un producto cartesiano}

Si se escriben varias tablas sin indicar cómo se relacionan, cada fila de una se combina con cada
fila de la otra. Ese producto cartesiano suele producir datos incoherentes.

\begin{lstlisting}
-- Incorrecto para relacionar empleados y cargos:
SELECT employee_id, job_title
FROM oehr_employees, oehr_jobs;
\end{lstlisting}

\subsection{INNER JOIN con ON}

Un \campo{INNER JOIN} devuelve coincidencias entre ambas tablas.

\begin{lstlisting}
SELECT e.employee_id,
       e.first_name || ' ' || e.last_name AS empleado,
       j.job_title,
       d.department_name
FROM oehr_employees e
JOIN oehr_jobs j
  ON e.job_id = j.job_id
JOIN oehr_departments d
  ON e.department_id = d.department_id;
\end{lstlisting}

Los alias \campo{e}, \campo{j} y \campo{d} permiten calificar columnas y eliminan ambigüedades.
Nunca escribas \campo{ON(job\_id = job\_id)}: parece comparar la columna consigo misma y, si el nombre
existe en ambas tablas, Oracle informa ambigüedad.

\subsection{USING y NATURAL JOIN}

\campo{USING(columna)} es posible cuando la columna de unión tiene el mismo nombre en ambas tablas.
Dentro de \campo{USING} no se califica la columna con alias.

\begin{lstlisting}
SELECT e.employee_id, d.department_name
FROM oehr_employees e
JOIN oehr_departments d
USING (department_id);
\end{lstlisting}

\campo{NATURAL JOIN} une automáticamente todas las columnas de igual nombre. Puede ser breve, pero
es menos explícito y puede cambiar de significado si el esquema incorpora otra columna coincidente.
Para la prueba, \campo{JOIN ... ON} suele mostrar mejor tu razonamiento.

\subsection{LEFT OUTER JOIN}

Una unión interna elimina al empleado cuyo \campo{DEPARTMENT\_ID} sea nulo o no encuentre
coincidencia. Si el requerimiento dice ``aunque no tenga departamento'', conserva la tabla de
empleados con \campo{LEFT JOIN}.

\begin{lstlisting}
SELECT e.employee_id,
       e.last_name,
       d.department_name
FROM oehr_employees e
LEFT JOIN oehr_departments d
  ON e.department_id = d.department_id;
\end{lstlisting}

La tabla escrita a la izquierda es la que se preserva. \campo{RIGHT OUTER JOIN} preserva la derecha;
\campo{FULL OUTER JOIN} preserva filas no coincidentes de ambas.

\subsection{Self join para empleado y jefe}

La misma tabla cumple dos roles. Se le asignan dos alias diferentes: \campo{e} es el empleado y
\campo{jefe} es la fila del jefe.

\begin{lstlisting}
SELECT e.employee_id,
       e.first_name || ' ' || e.last_name AS empleado,
       jefe.first_name || ' ' || jefe.last_name AS jefe
FROM oehr_employees e
LEFT JOIN oehr_employees jefe
  ON e.manager_id = jefe.employee_id;
\end{lstlisting}

Se usa \campo{LEFT JOIN} para no perder al empleado superior que no tiene jefe. Si el enunciado sólo
pide empleados que sí tienen jefe, un \campo{JOIN} interno es suficiente.

\subsection{Cadena de uniones}

Para mostrar la región en que trabaja el empleado, sigue cada clave foránea.

\begin{lstlisting}
SELECT e.employee_id,
       e.last_name,
       r.region_name
FROM oehr_employees e
JOIN oehr_departments d ON e.department_id = d.department_id
JOIN oehr_locations l   ON d.location_id   = l.location_id
JOIN oehr_countries c   ON l.country_id    = c.country_id
JOIN oehr_regions r     ON c.region_id     = r.region_id;
\end{lstlisting}

\subsection{Taller guiado: construir una unión desde cero}

Imagina el requerimiento ``mostrar todos los trabajadores, el nombre de su departamento, la ciudad
del departamento y el nombre de su cargo; deben aparecer incluso quienes no tengan departamento''.
No conviene intentar recordar una consulta completa. Se construye por capas. Primero fija la unidad
de salida: una fila por trabajador, de modo que \tabla{OEHR_EMPLOYEES e} debe iniciar el
\campo{FROM}. Segundo, marca dónde vive cada dato: nombre del departamento en
\tabla{OEHR_DEPARTMENTS d}, ciudad en \tabla{OEHR_LOCATIONS l} y título del cargo en
\tabla{OEHR_JOBS j}. Tercero, dibuja los enlaces: \campo{e.department_id} llega a
\campo{d.department_id}; desde esa fila, \campo{d.location_id} llega a \campo{l.location_id}; y
\campo{e.job_id} llega directamente a \campo{j.job_id}. Cuarto, decide qué relaciones son opcionales.
El enunciado obliga a conservar empleados sin departamento, por lo que el primer enlace departamental
es \campo{LEFT JOIN}. La ubicación depende del departamento: si éste falta, también debe poder faltar
la ubicación sin eliminar al empleado, por lo que el enlace hacia \campo{l} también es
\campo{LEFT JOIN}. En cambio, \campo{JOB_ID} es obligatorio según la definición de empleados; un
\campo{JOIN} interno con cargos es razonable. Quinto, escribe el \campo{SELECT} calificando cada
columna. Si más tarde se pide sólo empleados de Santiago, existe una decisión semántica importante:
\campo{WHERE l.city='Santiago'} eliminará las filas cuya ciudad sea nula y, en la práctica, perderás
los empleados sin departamento. Si quieres conservarlos, el requisito debe formularse de otra forma,
por ejemplo permitiendo explícitamente \campo{l.city IS NULL}. Esta interacción muestra que elegir
\campo{LEFT JOIN} no garantiza por sí sola la preservación: un filtro posterior sobre la tabla opcional
puede volver a descartarla. Finalmente, revisa la cantidad esperada de filas. Cada empleado tiene como
máximo un departamento, una ubicación y un cargo, así que esta cadena no debería multiplicar a un
empleado. Si apareciera repetido, revisa las claves antes de agregar \campo{DISTINCT}.

\begin{lstlisting}
SELECT e.employee_id,
       e.first_name || ' ' || e.last_name AS empleado,
       d.department_name,
       l.city,
       j.job_title
FROM oehr_employees e
LEFT JOIN oehr_departments d
  ON e.department_id = d.department_id
LEFT JOIN oehr_locations l
  ON d.location_id = l.location_id
JOIN oehr_jobs j
  ON e.job_id = j.job_id
ORDER BY d.department_name, e.last_name;
\end{lstlisting}

\begin{clave}
La ``X manera'' de juntar tablas no se adivina: se determina respondiendo tres preguntas. ¿Cuál es la
tabla que define una fila del resultado? ¿Qué pares PK--FK representan el camino hacia los datos que
faltan? ¿Deben sobrevivir las filas sin coincidencia? Las respuestas determinan el orden de las tablas,
las condiciones \campo{ON} y si corresponde \campo{JOIN} o \campo{LEFT JOIN}.
\end{clave}

\subsection{Unión no equitativa y CROSS JOIN}

Una unión no equitativa usa un operador distinto de igualdad, por ejemplo para ubicar un sueldo entre
un mínimo y un máximo. Sólo puede practicarse si existe una tabla de rangos adecuada.
\campo{CROSS JOIN} expresa deliberadamente un producto cartesiano; no debe aparecer por accidente.

\subsection{Lectura activa del módulo 6}

\begin{enumerate}[label=\textbf{M6.\arabic*}]
  \item \nivel{Muy fácil.} ¿Qué cláusula indica la condición que conecta dos tablas?
  \item \nivel{Muy fácil.} ¿Para qué sirven los alias de tabla?
  \item \nivel{Muy fácil.} ¿Qué tipo de unión conserva empleados sin departamento?
  \item \nivel{Intermedia.} Muestra código, apellido y título del cargo de cada empleado.
  \item \nivel{Intermedia.} Muestra apellido del empleado, departamento y ciudad.
  \item \nivel{Difícil.} Muestra código, nombre completo, nombre del departamento, nombre completo
  del jefe del departamento, ciudad y título del cargo. Conserva al empleado aunque su departamento
  o el jefe del departamento sea nulo.
\end{enumerate}

\section{Módulo 7: subconsultas}

\begin{objetivo}
Usar el resultado de una consulta como criterio de otra consulta y decidir si corresponde
\campo{=}, \campo{IN} u otro operador.
\end{objetivo}

Una subconsulta resuelve una pregunta auxiliar cuyo resultado necesita la consulta principal. La
forma más segura de construirla es verbalizar dos preguntas separadas. Para ``mostrar todos los
empleados que trabajan en los departamentos de Whalen o King'', la pregunta interior es ``¿cuáles son
los códigos de departamento de Whalen o King?'' y devuelve una lista de
\campo{DEPARTMENT_ID}. La pregunta exterior es ``¿qué empleados tienen un código contenido en esa
lista?''. Como puede haber más de un código, se usa \campo{IN}; emplear \campo{=} supondría una sola
fila y produciría un error si aparecen varias. El hecho de que una consulta esté escrita dentro de
otra no cambia las reglas ya aprendidas: la interna posee su propio \campo{SELECT}, \campo{FROM} y
\campo{WHERE}, y puede incluir agrupamiento. Conviene ejecutarla mentalmente o por separado, comprobar
qué columna y cuántas filas produce, y recién entonces elegir el operador exterior. Una subconsulta
que calcula \campo{AVG(salary)} sin agrupar garantiza un solo valor y puede compararse con
\campo{>}; una que selecciona cargos de un departamento normalmente produce varios valores y pide
\campo{IN}. Luego debes distinguir subconsulta de unión. Una unión aporta columnas de otra tabla a
cada fila de salida; una subconsulta usada en \campo{WHERE} aporta un criterio para decidir qué filas
pasan. Ambas pueden convivir: la consulta exterior une departamentos para mostrar su nombre, mientras
la interior descubre los códigos permitidos. Presta atención a los nulos con \campo{NOT IN}: si la
lista interior contiene un nulo, las comparaciones pueden quedar desconocidas. En ejercicios
introductorios puedes filtrar \campo{IS NOT NULL}; más adelante \campo{NOT EXISTS} ofrece otra forma
de expresar ausencia. El objetivo no es anidar por complejidad, sino dividir un problema grande en
preguntas pequeñas cuyos resultados encajan por tipo y cardinalidad.

Una subconsulta es una consulta encerrada entre paréntesis. Primero se razona la consulta interior y
luego se usa su resultado en la exterior.

\subsection{Subconsulta de una fila}

Si la consulta interior garantiza un solo valor, puede compararse con \campo{=}, \campo{>}, etc.

\begin{lstlisting}
SELECT employee_id, last_name, salary
FROM oehr_employees
WHERE salary > (
  SELECT AVG(salary)
  FROM oehr_employees
);
\end{lstlisting}

La consulta interior produce un promedio; la exterior conserva empleados sobre ese promedio.

\subsection{Subconsulta de varias filas con IN}

Si la subconsulta puede devolver varios departamentos, \campo{=} no sirve; se necesita \campo{IN}.

\begin{lstlisting}
SELECT e.employee_id,
       e.first_name || ' ' || e.last_name AS empleado,
       d.department_name
FROM oehr_employees e
JOIN oehr_departments d
  ON e.department_id = d.department_id
WHERE e.department_id IN (
  SELECT department_id
  FROM oehr_employees
  WHERE UPPER(last_name) IN ('WHALEN', 'KING')
);
\end{lstlisting}

Éste es el patrón de la quinta consulta visible en la pauta. La subconsulta responde ``¿en qué
departamentos trabajan Whalen o King?''; la consulta externa obtiene todos los empleados de esos
departamentos.

\subsection{Subconsultas y valores nulos}

\campo{NOT IN} puede dar resultados inesperados si la subconsulta devuelve algún nulo, porque la
comparación pasa a ser desconocida. Cuando exista esa posibilidad, filtra los nulos en la subconsulta
o estudia \campo{NOT EXISTS}. Para el patrón confirmado de la pauta se usa \campo{IN}, por lo que el
problema no es central, pero conviene reconocerlo.

\subsection{Lectura activa del módulo 7}

\begin{enumerate}[label=\textbf{M7.\arabic*}]
  \item \nivel{Muy fácil.} ¿Entre qué signos se escribe una subconsulta?
  \item \nivel{Muy fácil.} Si la subconsulta devuelve varios valores, ¿qué operador permite comparar
  contra todos ellos?
  \item \nivel{Muy fácil.} ¿Qué consulta debes comprender primero: la interna o la externa?
  \item \nivel{Intermedia.} Muestra empleados que ganan más que el sueldo promedio del departamento
  90.
  \item \nivel{Intermedia.} Muestra empleados cuyo cargo sea uno de los cargos que existen en el
  departamento 50.
  \item \nivel{Difícil.} Muestra código, nombre y departamento de quienes trabajan en un departamento
  que contiene al menos un empleado cuyo sueldo supera 12000. No escribas manualmente los códigos de
  departamento.
\end{enumerate}

\ifmostrarSoluciones
\section{Pauta razonada de los seis requerimientos}

Esta pauta se alinea con los seis requerimientos del documento de evaluación. No debe memorizarse
literalmente: primero resuelve cada ejercicio, luego compara la estructura y explica con tus propias
palabras por qué se utilizó cada cláusula.

\subsection{Ejercicio 1: texto y patrón}

\textbf{Habilidades:} concatenación, \campo{UPPER}, \campo{LIKE}, \campo{OR}.

\begin{lstlisting}
SELECT 'El codigo del cargo ' || job_title ||
       ' es: ' || job_id AS "CARGO",
       min_salary AS "SUELDO MINIMO"
FROM oehr_jobs
WHERE UPPER(job_title) LIKE 'S%'
   OR UPPER(job_id)    LIKE 'S%';
\end{lstlisting}

Comprueba que la frase, el sueldo mínimo y ambos criterios de búsqueda estén presentes. El nombre
definido en el esquema es \tabla{OEHR\_JOBS}.

\subsection{Ejercicio 2: resumen por cargo}

\textbf{Habilidades:} \campo{COUNT}, \campo{AVG}, \campo{MIN}, \campo{MAX},
\campo{GROUP BY} y \campo{HAVING}.

\begin{lstlisting}
SELECT job_id AS "OFICIO",
       COUNT(employee_id) AS "CANTIDAD EMPLEADOS",
       ROUND(AVG(salary), 2) AS "PROMEDIO DE SUELDOS",
       MAX(salary) - MIN(salary) AS "DIFERENCIA SUELDOS"
FROM oehr_employees
GROUP BY job_id
HAVING MAX(salary) - MIN(salary) <> 0;
\end{lstlisting}

La condición agregada pertenece a \campo{HAVING}, porque se evalúa después de formar cada grupo. La
diferencia también debe aparecer en la salida, ya que el enunciado la solicita explícitamente.

\subsection{Ejercicio 3: empleado, departamento y jefe}

\textbf{Habilidades:} unión normal, auto-unión y alias de tabla. El enunciado pide al jefe del
departamento, por lo que el segundo enlace nace en \campo{d.manager\_id}; no en
\campo{e.manager\_id}, que representa al jefe directo del empleado.

\begin{lstlisting}
SELECT e.employee_id,
       e.first_name || ' ' || e.last_name AS empleado,
       d.department_name,
       jefe.first_name || ' ' || jefe.last_name AS jefe
FROM oehr_employees e
JOIN oehr_departments d
  ON e.department_id = d.department_id
JOIN oehr_employees jefe
  ON d.manager_id = jefe.employee_id;
\end{lstlisting}

Si una variante exigiera conservar empleados sin departamento o departamentos sin jefe, se usarían
\campo{LEFT JOIN}. En la respuesta de referencia se emplea \campo{JOIN}, porque se solicitan datos
asociados de las tres filas.

\subsection{Ejercicio 4: promedio del mínimo salarial por cargo}

El fragmento de pauta es ambiguo: \tabla{OEHR\_JOBS} posee una sola fila por \campo{JOB\_ID}, de modo
que promediar \campo{MIN\_SALARY} por ese mismo código no agrega información. Una reconstrucción
sintácticamente válida es:

\begin{lstlisting}
SELECT job_id,
       AVG(min_salary) AS promedio_minimo
FROM oehr_jobs
GROUP BY job_id
HAVING AVG(min_salary) > 3000;
\end{lstlisting}

Pedagógicamente, lo importante es reconocer que una condición sobre \campo{AVG} se escribe en
\campo{HAVING}. En una prueba real se debe leer el enunciado completo antes de decidir si la tabla o
el agrupamiento son realmente esos.

\subsection{Ejercicio 5: departamentos de Whalen o King}

\textbf{Habilidades:} \campo{JOIN}, concatenación y subconsulta multirregistro con \campo{IN}.

\begin{lstlisting}
SELECT e.employee_id,
       e.first_name || ' ' || e.last_name AS empleado,
       d.department_name
FROM oehr_employees e
JOIN oehr_departments d
  ON e.department_id = d.department_id
WHERE e.department_id IN (
  SELECT department_id
  FROM oehr_employees
  WHERE UPPER(last_name) IN ('WHALEN', 'KING')
)
ORDER BY d.department_name, e.last_name;
\end{lstlisting}

La subconsulta puede devolver más de un departamento: King no identifica necesariamente una única
fila. Por esa razón se compara con \campo{IN} y no con el operador \campo{=}.

\subsection{Ejercicio 6: comisión, rango, correo y departamento de Kochhar}

\textbf{Habilidades:} concatenación, \campo{SUBSTR}, nulos, \campo{NOT BETWEEN}, conectores lógicos
y subconsulta.

\begin{lstlisting}
SELECT e.first_name || ' ' || e.last_name AS "NOMBRE EMPLEADO",
       SUBSTR(e.first_name, 1, 1) ||
       SUBSTR(e.last_name, 2, 1) ||
       e.job_id || '@pruebasql.cl' AS "NUEVO CORREO",
       e.commission_pct AS "COMISION"
FROM oehr_employees e
WHERE e.commission_pct IS NOT NULL
  AND e.salary NOT BETWEEN 5000 AND 10000
  AND e.department_id IN (
    SELECT department_id
    FROM oehr_employees
    WHERE UPPER(first_name) = 'NEENA'
      AND UPPER(last_name) = 'KOCHHAR'
  );
\end{lstlisting}

La frase ``tengan una comisión asociada'' se interpreta como
\campo{commission\_pct IS NOT NULL}. ``Fuera del rango'' se expresa con
\campo{NOT BETWEEN}; los extremos 5000 y 10000 quedan fuera del resultado porque
\campo{BETWEEN} es inclusivo.
\fi

\section{Lista de errores que debes poder diagnosticar}

\begin{longtable}{p{0.42\textwidth} p{0.50\textwidth}}
\toprule
\textbf{Error} & \textbf{Corrección razonada} \\
\midrule
\endhead
\campo{OHER\_JOBS} & El esquema define \campo{OEHR\_JOBS}. \\
\campo{DEPARMENT\_ID} & La columna es \campo{DEPARTMENT\_ID}. \\
Alias con comillas simples & Usa alias sin comillas o con comillas dobles. \\
Falta una coma en \campo{SELECT} & Cada expresión seleccionada se separa con coma. \\
Punto y coma antes de terminar & Escribe un único punto y coma al final. \\
\campo{WHERE MAX(salary)>...} & Una condición agregada corresponde a \campo{HAVING}. \\
Columna normal junto a \campo{AVG} sin agrupar & Agrega la columna a \campo{GROUP BY} o quítala del
\campo{SELECT}. \\
\campo{commission\_pct = NULL} & Usa \campo{IS NULL}. \\
\campo{ON(job\_id=job\_id)} & Califica ambos lados con alias de tabla. \\
Varias tablas sin unión & Escribe cada relación mediante \campo{JOIN ... ON}. \\
\campo{CONCAT(a,b,c)} & Usa \campo{a || b || c} o anida \campo{CONCAT}. \\
Fecha como texto ambiguo & Usa \campo{TO\_DATE(texto,formato)}. \\
\campo{MOD(54/3)} & La función recibe dos argumentos: \campo{MOD(54,3)}. \\
Cadena de \campo{<>} unida con \campo{OR} & Usa \campo{NOT IN (...)} o une correctamente con
\campo{AND}. \\
\bottomrule
\end{longtable}

\section{Pauta de autoevaluación alineada con la rúbrica}

Para cada fila marca los criterios comprobados y luego asigna un nivel: Destacado = 5,
Competente = 4, Básico = 3 y En desarrollo = 1. Los criterios internos orientan la corrección, pero
cada fila aporta como máximo 5 puntos. La suma máxima es 35.

\begin{longtable}{p{0.13\textwidth} p{0.67\textwidth} p{0.10\textwidth}}
\toprule
\textbf{Ítem} & \textbf{Criterios observables} & \textbf{Nivel} \\
\midrule
\endhead
Formalidad & \casilla{} archivo \campo{.sql}; \casilla{} nombre del archivo correcto;
\casilla{} encabezado con datos; \casilla{} un comentario por requerimiento; \casilla{} código
ordenado y entregable. & /5 \\
R1 & \casilla{} muestra frase y sueldo mínimo; \casilla{} concatena; \casilla{} usa
\campo{UPPER}; \casilla{} filtra con \campo{LIKE 'S\%'}; \casilla{} conecta ambos criterios con
\campo{OR}. & /5 \\
R2 & \casilla{} muestra las cuatro métricas; \casilla{} usa \campo{COUNT} y \campo{AVG};
\casilla{} calcula \campo{MAX-MIN}; \casilla{} agrupa por oficio; \casilla{} filtra la diferencia
con \campo{HAVING}. & /5 \\
R3 & \casilla{} muestra los cuatro datos; \casilla{} usa dos \campo{JOIN}; \casilla{} escribe cada
\campo{ON}; \casilla{} usa alias de tabla; \casilla{} obtiene al jefe desde
\campo{d.manager\_id}. & /5 \\
R4 & \casilla{} muestra el código de oficio; \casilla{} calcula \campo{AVG(min\_salary)};
\casilla{} agrupa por oficio; \casilla{} usa \campo{HAVING}; \casilla{} compara el promedio con
3000. & /5 \\
R5 & \casilla{} muestra código, nombre y departamento; \casilla{} concatena nombre y apellido;
\casilla{} une empleados y departamentos; \casilla{} la subconsulta busca Whalen o King;
\casilla{} compara con \campo{IN}. & /5 \\
R6 & \casilla{} construye el correo pedido; \casilla{} exige comisión con \campo{IS NOT NULL};
\casilla{} usa \campo{NOT BETWEEN}; \casilla{} combina condiciones con \campo{AND};
\casilla{} obtiene el departamento de Neena Kochhar mediante subconsulta. & /5 \\
\midrule
\textbf{Total} & Escribe el error que debes corregir primero en el próximo intento:
\rule{7cm}{0.4pt} & \textbf{/35} \\
\bottomrule
\end{longtable}

\begin{pausa}
No te asignes puntaje sólo porque la consulta ``se parece'' a la pauta. Comprueba que ejecuta, que
las columnas corresponden al enunciado y que las filas obtenidas cumplen todas las condiciones. Si
no puedes ejecutarla, recorre en papel las seis decisiones y revisa cada relación de claves.
\end{pausa}

\section{Simulacro acumulativo}

Resuelve sin mirar el solucionario. Antes de escribir SQL, anota las seis decisiones del método
inicial.

\subsection*{Parte A: muy fácil}

\begin{enumerate}[label=\textbf{S.\arabic*}]
  \item Muestra código, apellido, sueldo y departamento de todos los empleados.
  \item Muestra los códigos de cargo distintos presentes en empleados.
  \item Muestra los empleados que no reciben comisión.
\end{enumerate}

\subsection*{Parte B: intermedia}

\begin{enumerate}[label=\textbf{S.\arabic*},resume]
  \item Muestra código, nombre completo en mayúsculas, correo terminado en \campo{@BD.UCM.CL} y fecha
  de contratación con formato \campo{DD-MM-YYYY}, para departamentos 50 y 80. Ordena desde la
  contratación más reciente.
  \item Muestra nombre de departamento, cantidad de empleados, sueldo mínimo, máximo y promedio.
  Conserva departamentos cuyo promedio sea mayor que 7000.
\end{enumerate}

\subsection*{Parte C: difícil}

\begin{enumerate}[label=\textbf{S.\arabic*},resume]
  \item Muestra el código y nombre del empleado, título de su cargo, nombre de su departamento,
  ciudad, región y nombre de su jefe directo. Conserva empleados sin jefe o sin departamento. La
  salida debe estar ordenada por región, departamento y apellido.
\end{enumerate}

\ifmostrarSoluciones
\newpage
\section{Solucionario razonado}

No uses este apartado como lectura pasiva. Para cada respuesta, señala con un lápiz qué parte
corresponde a salida, origen, relación, filtro de fila, agrupamiento, filtro de grupo y orden.

\subsection{Soluciones del módulo 1}

\begin{enumerate}[label=\textbf{M1.\arabic*}]
  \item \campo{SELECT}.
  \item El asterisco \campo{*}.
  \item \campo{DISTINCT}.
  \item
\begin{lstlisting}
SELECT employee_id AS "CODIGO",
       last_name AS "APELLIDO",
       salary * 12 AS "ANUAL"
FROM oehr_employees;
\end{lstlisting}
  \item La primera suma sólo 1200 al sueldo por la precedencia de la multiplicación. La segunda suma
  100 al sueldo mensual y luego anualiza todo.
  \item
\begin{lstlisting}
SELECT UPPER(last_name) || ' gana ' ||
       TO_CHAR(salary * 12, 'FM999999990') ||
       ' al anio' AS reporte
FROM oehr_employees;
\end{lstlisting}
\end{enumerate}

\subsection{Soluciones del módulo 2}

\begin{enumerate}[label=\textbf{M2.\arabic*}]
  \item \campo{IS NULL}.
  \item Exactamente un carácter.
  \item \campo{DESC}.
  \item
\begin{lstlisting}
SELECT employee_id, last_name, department_id, salary
FROM oehr_employees
WHERE department_id IN (50, 80)
  AND salary BETWEEN 4000 AND 9000;
\end{lstlisting}
  \item \campo{WHERE LOWER(last\_name) LIKE '\_a\%s'}.
  \item
\begin{lstlisting}
SELECT employee_id, last_name, department_id, salary
FROM oehr_employees
WHERE department_id IN (50, 80, 90)
  AND commission_pct IS NULL
ORDER BY department_id ASC, salary DESC;
\end{lstlisting}
\end{enumerate}

\subsection{Soluciones del módulo 3}

\begin{enumerate}[label=\textbf{M3.\arabic*}]
  \item \campo{UPPER}.
  \item \campo{ROUND} redondea; \campo{TRUNC} elimina las posiciones sobrantes sin redondear.
  \item Para evaluar una expresión una vez cuando no depende de una tabla de datos del usuario.
  \item
\begin{lstlisting}
SELECT employee_id,
       LOWER(SUBSTR(first_name, 1, 2) || '.' ||
             last_name || '@ucm.cl') AS correo
FROM oehr_employees;
\end{lstlisting}
  \item
\begin{lstlisting}
SELECT employee_id,
       TRUNC(MONTHS_BETWEEN(SYSDATE, hire_date) / 12) AS anios
FROM oehr_employees;
\end{lstlisting}
  \item
\begin{lstlisting}
SELECT employee_id,
       UPPER(SUBSTR(last_name, -1, 1) ||
             SUBSTR(first_name, 1, 1)) ||
       TO_CHAR(LENGTH(first_name || last_name)) AS identificador
FROM oehr_employees;
\end{lstlisting}
\end{enumerate}

\subsection{Soluciones del módulo 4}

\begin{enumerate}[label=\textbf{M4.\arabic*}]
  \item \campo{TO\_CHAR}.
  \item Cero.
  \item \campo{COALESCE}.
  \item
\begin{lstlisting}
SELECT employee_id,
       TO_CHAR(salary, 'FM999999990.00') AS sueldo,
       TO_CHAR(hire_date, 'YYYY-MM-DD') AS fecha
FROM oehr_employees;
\end{lstlisting}
  \item
\begin{lstlisting}
SELECT employee_id,
       CASE
         WHEN commission_pct IS NULL THEN 'SIN COMISION'
         ELSE 'CON COMISION'
       END AS estado
FROM oehr_employees;
\end{lstlisting}
  \item
\begin{lstlisting}
SELECT first_name || ' ' || last_name AS empleado,
       TO_CHAR(salary, 'FM$999,999,990.00') AS sueldo,
       CASE
         WHEN commission_pct IS NULL THEN 'SIN COMISION'
         ELSE TO_CHAR(salary * commission_pct,
                      'FM$999,999,990.00')
       END AS monto_comision
FROM oehr_employees;
\end{lstlisting}
\end{enumerate}

\subsection{Soluciones del módulo 5}

\begin{enumerate}[label=\textbf{M5.\arabic*}]
  \item \campo{AVG}.
  \item \campo{COUNT(*)} cuenta filas; \campo{COUNT(commission\_pct)} sólo cuenta filas donde esa
  columna no es nula.
  \item \campo{HAVING}.
  \item
\begin{lstlisting}
SELECT job_id, COUNT(*) AS cantidad
FROM oehr_employees
GROUP BY job_id;
\end{lstlisting}
  \item
\begin{lstlisting}
SELECT department_id, ROUND(AVG(salary), 2) AS promedio
FROM oehr_employees
GROUP BY department_id
HAVING AVG(salary) > 8000;
\end{lstlisting}
  \item
\begin{lstlisting}
SELECT department_id,
       job_id,
       COUNT(*) AS cantidad,
       ROUND(AVG(salary), 2) AS promedio
FROM oehr_employees
WHERE department_id IN (50, 80, 90)
GROUP BY department_id, job_id
HAVING COUNT(*) >= 2
ORDER BY promedio DESC;
\end{lstlisting}
\end{enumerate}

\subsection{Soluciones del módulo 6}

\begin{enumerate}[label=\textbf{M6.\arabic*}]
  \item \campo{ON}; también puede utilizarse \campo{USING} cuando corresponda.
  \item Para representar el rol de cada tabla, acortar sus nombres y evitar columnas ambiguas.
  \item \campo{LEFT JOIN} desde empleados hacia departamentos.
  \item
\begin{lstlisting}
SELECT e.employee_id, e.last_name, j.job_title
FROM oehr_employees e
JOIN oehr_jobs j ON e.job_id = j.job_id;
\end{lstlisting}
  \item
\begin{lstlisting}
SELECT e.last_name, d.department_name, l.city
FROM oehr_employees e
JOIN oehr_departments d ON e.department_id = d.department_id
JOIN oehr_locations l ON d.location_id = l.location_id;
\end{lstlisting}
  \item El ``jefe del departamento'' se obtiene desde el identificador de jefe de la tabla
  departamentos (alias \campo{d}), no necesariamente desde el jefe directo del empleado.
\begin{lstlisting}
SELECT e.employee_id,
       e.first_name || ' ' || e.last_name AS empleado,
       d.department_name,
       jd.first_name || ' ' || jd.last_name AS jefe_departamento,
       l.city,
       c.job_title
FROM oehr_employees e
JOIN oehr_jobs c
  ON e.job_id = c.job_id
LEFT JOIN oehr_departments d
  ON e.department_id = d.department_id
LEFT JOIN oehr_employees jd
  ON d.manager_id = jd.employee_id
LEFT JOIN oehr_locations l
  ON d.location_id = l.location_id;
\end{lstlisting}
\end{enumerate}

\subsection{Soluciones del módulo 7}

\begin{enumerate}[label=\textbf{M7.\arabic*}]
  \item Entre paréntesis.
  \item \campo{IN}.
  \item La interna, porque su resultado alimenta la condición externa.
  \item
\begin{lstlisting}
SELECT employee_id, last_name, salary
FROM oehr_employees
WHERE salary > (
  SELECT AVG(salary)
  FROM oehr_employees
  WHERE department_id = 90
);
\end{lstlisting}
  \item
\begin{lstlisting}
SELECT employee_id, last_name, job_id
FROM oehr_employees
WHERE job_id IN (
  SELECT job_id
  FROM oehr_employees
  WHERE department_id = 50
);
\end{lstlisting}
  \item
\begin{lstlisting}
SELECT e.employee_id,
       e.first_name || ' ' || e.last_name AS empleado,
       d.department_name
FROM oehr_employees e
JOIN oehr_departments d
  ON e.department_id = d.department_id
WHERE e.department_id IN (
  SELECT department_id
  FROM oehr_employees
  WHERE salary > 12000
    AND department_id IS NOT NULL
);
\end{lstlisting}
\end{enumerate}

\subsection{Soluciones del simulacro}

\begin{enumerate}[label=\textbf{S.\arabic*}]
  \item
\begin{lstlisting}
SELECT employee_id, last_name, salary, department_id
FROM oehr_employees;
\end{lstlisting}
  \item
\begin{lstlisting}
SELECT DISTINCT job_id
FROM oehr_employees;
\end{lstlisting}
  \item
\begin{lstlisting}
SELECT employee_id, last_name
FROM oehr_employees
WHERE commission_pct IS NULL;
\end{lstlisting}
  \item
\begin{lstlisting}
SELECT employee_id,
       UPPER(first_name || ' ' || last_name) AS nombre,
       email || '@BD.UCM.CL' AS correo,
       TO_CHAR(hire_date, 'DD-MM-YYYY') AS fecha_contrato
FROM oehr_employees
WHERE department_id IN (50, 80)
ORDER BY hire_date DESC;
\end{lstlisting}
  \item Se usa una unión para mostrar el nombre y se agrupa por identificador y nombre para no
  mezclar departamentos distintos que, hipotéticamente, tuvieran el mismo nombre.
\begin{lstlisting}
SELECT d.department_id,
       d.department_name,
       COUNT(e.employee_id) AS cantidad,
       MIN(e.salary) AS minimo,
       MAX(e.salary) AS maximo,
       ROUND(AVG(e.salary), 2) AS promedio
FROM oehr_departments d
JOIN oehr_employees e
  ON d.department_id = e.department_id
GROUP BY d.department_id, d.department_name
HAVING AVG(e.salary) > 7000;
\end{lstlisting}
  \item La preservación de empleados exige que todas las tablas opcionales cuelguen mediante
  \campo{LEFT JOIN} desde la cadena preservada.
\begin{lstlisting}
SELECT e.employee_id,
       e.first_name || ' ' || e.last_name AS empleado,
       j.job_title,
       d.department_name,
       l.city,
       r.region_name,
       jefe.first_name || ' ' || jefe.last_name AS jefe
FROM oehr_employees e
JOIN oehr_jobs j
  ON e.job_id = j.job_id
LEFT JOIN oehr_departments d
  ON e.department_id = d.department_id
LEFT JOIN oehr_locations l
  ON d.location_id = l.location_id
LEFT JOIN oehr_countries c
  ON l.country_id = c.country_id
LEFT JOIN oehr_regions r
  ON c.region_id = r.region_id
LEFT JOIN oehr_employees jefe
  ON e.manager_id = jefe.employee_id
ORDER BY r.region_name, d.department_name, e.last_name;
\end{lstlisting}
\end{enumerate}

\fi
\section{Hoja de repaso final}

Antes de entregar una consulta, verifica:

\begin{enumerate}
  \item ¿Escribí exactamente \tabla{OEHR\_EMPLOYEES}, \tabla{OEHR\_DEPARTMENTS} y
  \tabla{OEHR\_JOBS}?
  \item ¿Cada columna pertenece realmente a la tabla cuyo alias utilicé?
  \item ¿Separé todas las expresiones del \campo{SELECT} con comas?
  \item ¿Las cadenas tienen comillas simples y los alias especiales comillas dobles?
  \item ¿Comparé los nulos con \campo{IS NULL} o \campo{IS NOT NULL}?
  \item ¿Normalicé mayúsculas y minúsculas en búsquedas textuales?
  \item ¿Usé \campo{WHERE} para filas y \campo{HAVING} para resultados agregados?
  \item ¿Incluí en \campo{GROUP BY} todas las columnas no agregadas del resultado?
  \item ¿Cada \campo{JOIN} tiene una relación de clave correcta y calificada con alias?
  \item ¿Elegí \campo{LEFT JOIN} cuando debían conservarse filas sin coincidencia?
  \item ¿Usé \campo{IN} si una subconsulta puede devolver varias filas?
  \item ¿Dejé un único punto y coma al final?
\end{enumerate}

\begin{clave}
La mejor forma de estudiar no es memorizar el orden visual de un ejemplo. Toma un requerimiento,
subraya sustantivos para descubrir columnas y tablas, encierra las restricciones de filas, marca las
palabras ``por cada'' para descubrir agrupamientos y busca ``promedio'', ``cantidad'', ``máximo'',
``mínimo'' o ``total'' para descubrir funciones de grupo.
\end{clave}

\section{Documentación oficial utilizada}

Las siguientes páginas pertenecen a \emph{Oracle Database 21c SQL Language Reference}. Se usaron
para comprobar la sintaxis y complementar los scripts del curso:

\begin{itemize}
  \item \href{https://docs.oracle.com/en/database/oracle/oracle-database/21/sqlrf/SELECT.html}
  {SELECT y sus cláusulas};
  \item \href{https://docs.oracle.com/en/database/oracle/oracle-database/21/sqlrf/Conditions.html}
  {condiciones SQL};
  \item \href{https://docs.oracle.com/en/database/oracle/oracle-database/21/sqlrf/Pattern-matching-Conditions.html}
  {condiciones de patrones LIKE};
  \item \href{https://docs.oracle.com/en/database/oracle/oracle-database/21/sqlrf/Null-Conditions.html}
  {condiciones de valores nulos};
  \item \href{https://docs.oracle.com/en/database/oracle/oracle-database/21/sqlrf/Single-Row-Functions.html}
  {funciones de una sola fila};
  \item \href{https://docs.oracle.com/en/database/oracle/oracle-database/21/sqlrf/Aggregate-Functions.html}
  {funciones de agregación};
  \item \href{https://docs.oracle.com/en/database/oracle/oracle-database/21/sqlrf/Joins.html}
  {uniones};
  \item \href{https://docs.oracle.com/en/database/oracle/oracle-database/21/sqlrf/Scalar-Subquery-Expressions.html}
  {expresiones de subconsulta escalar};
  \item \href{https://docs.oracle.com/en/database/oracle/oracle-database/21/sqlrf/TO_CHAR-datetime.html}
  {TO\_CHAR para fechas};
  \item \href{https://docs.oracle.com/en/database/oracle/oracle-database/21/sqlrf/TO_DATE.html}
  {TO\_DATE};
  \item \href{https://docs.oracle.com/en/database/oracle/oracle-database/21/sqlrf/NVL.html}
  {NVL} y
  \href{https://docs.oracle.com/en/database/oracle/oracle-database/21/sqlrf/NVL2.html}
  {NVL2};
  \item \href{https://docs.oracle.com/en/database/oracle/oracle-database/21/sqlrf/CASE-Expressions.html}
  {expresiones CASE} y
  \href{https://docs.oracle.com/en/database/oracle/oracle-database/21/sqlrf/DECODE.html}
  {DECODE};
  \item \href{https://docs.oracle.com/en/database/oracle/oracle-database/21/sqlrf/Datetime-Functions.html}
  {funciones de fecha y hora}.
\end{itemize}

\paragraph{Nota de alcance.}
La documentación oficial describe muchas capacidades que exceden esta evaluación. La selección de
temas de esta guía se restringe a lo enseñado en U1 L1--L6 y a las habilidades visibles en la pauta.
No se incluyen procedimientos PL/SQL, triggers, control de concurrencia ni administración, porque no
aparecen en los materiales de esta primera prueba.

\end{document}
